%%%%%%%%%%%%%%%%%%%%%%%%%%%%%%%%%%%%%%%%%%%%%%%%%%%%%%%%%%%%%%%%%%%%%%%%%%%%%%%%
% Template for USENIX papers.
%
% History:
%
% - TEMPLATE for Usenix papers, specifically to meet requirements of
%   USENIX '05. originally a template for producing IEEE-format
%   articles using LaTeX. written by Matthew Ward, CS Department,
%   Worcester Polytechnic Institute. adapted by David Beazley for his
%   excellent SWIG paper in Proceedings, Tcl 96. turned into a
%   smartass generic template by De Clarke, with thanks to both the
%   above pioneers. Use at your own risk. Complaints to /dev/null.
%   Make it two column with no page numbering, default is 10 point.
%
% - Munged by Fred Douglis <douglis@research.att.com> 10/97 to
%   separate the .sty file from the LaTeX source template, so that
%   people can more easily include the .sty file into an existing
%   document. Also changed to more closely follow the style guidelines
%   as represented by the Word sample file.
%
% - Note that since 2010, USENIX does not require endnotes. If you
%   want foot of page notes, don't include the endnotes package in the
%   usepackage command, below.
% - This version uses the latex2e styles, not the very ancient 2.09
%   stuff.
%
% - Updated July 2018: Text block size changed from 6.5" to 7"
%
% - Updated Dec 2018 for ATC'19:
%
%   * Revised text to pass HotCRP's auto-formatting check, with
%     hotcrp.settings.submission_form.body_font_size=10pt, and
%     hotcrp.settings.submission_form.line_height=12pt
%
%   * Switched from \endnote-s to \footnote-s to match Usenix's policy.
%
%   * \section* => \begin{abstract} ... \end{abstract}
%
%   * Make template self-contained in terms of bibtex entires, to allow
%     this file to be compiled. (And changing refs style to 'plain'.)
%
%   * Make template self-contained in terms of figures, to
%     allow this file to be compiled. 
%
%   * Added packages for hyperref, embedding fonts, and improving
%     appearance.
%   
%   * Removed outdated text.
%
%%%%%%%%%%%%%%%%%%%%%%%%%%%%%%%%%%%%%%%%%%%%%%%%%%%%%%%%%%%%%%%%%%%%%%%%%%%%%%%%
\PassOptionsToPackage{table}{xcolor}
\documentclass[letterpaper,twocolumn,10pt]{article}
\usepackage{usenix}

% to be able to draw some self-contained figs
%%%%%%%%%%%%%%%%%%%%%%%%%%%%%%%%%%%%%%%%%%%%%%%%%%%%%%%%%%%%%%%%%%%%%%%%%%%%%%%%
% Mathematics
%%%%%%%%%%%%%%%%%%%%%%%%%%%%%%%%%%%%%%%%%%%%%%%%%%%%%%%%%%%%%%%%%%%%%%%%%%%%%%%%
\usepackage{bm}
\usepackage{amsfonts}
\usepackage{amsmath}
\usepackage{mathtools}
\usepackage{amssymb}
\usepackage{amsthm}

%%%%%%%%%%%%%%%%%%%%%%%%%%%%%%%%%%%%%%%%%%%%%%%%%%%%%%%%%%%%%%%%%%%%%%%%%%%%%%%%
% Algorithms
%%%%%%%%%%%%%%%%%%%%%%%%%%%%%%%%%%%%%%%%%%%%%%%%%%%%%%%%%%%%%%%%%%%%%%%%%%%%%%%%
\usepackage{algorithm}
\usepackage{algpseudocode}

%%%%%%%%%%%%%%%%%%%%%%%%%%%%%%%%%%%%%%%%%%%%%%%%%%%%%%%%%%%%%%%%%%%%%%%%%%%%%%%%
% Tables and figures
%%%%%%%%%%%%%%%%%%%%%%%%%%%%%%%%%%%%%%%%%%%%%%%%%%%%%%%%%%%%%%%%%%%%%%%%%%%%%%%%
\usepackage{booktabs}
\usepackage{multirow}
\usepackage{graphicx}
\usepackage{subcaption}
\usepackage{tikz}
\usepackage{colortbl}
\usetikzlibrary{arrows.meta, positioning, shapes.multipart, trees}
\usepackage{siunitx}
\usepackage{makecell}

%%%%%%%%%%%%%%%%%%%%%%%%%%%%%%%%%%%%%%%%%%%%%%%%%%%%%%%%%%%%%%%%%%%%%%%%%%%%%%%%
% Text utilities
%%%%%%%%%%%%%%%%%%%%%%%%%%%%%%%%%%%%%%%%%%%%%%%%%%%%%%%%%%%%%%%%%%%%%%%%%%%%%%%%
\usepackage[normalem]{ulem}

\usepackage{kotex}

%%%%%%%%%%%%%%%%%%%%%%%%%%%%%%%%%%%%%%%%%%%%%%%%%%%%%%%%%%%%%%%%%%%%%%%%%%%%%%%%
% References
%%%%%%%%%%%%%%%%%%%%%%%%%%%%%%%%%%%%%%%%%%%%%%%%%%%%%%%%%%%%%%%%%%%%%%%%%%%%%%%%

% hyperref는 usenix.sty가 이미 불러온다.
% cleveref는 hyperref 및 다른 환경들 뒤에 불러오는 것이 좋다.
\usepackage{cleveref}

%%%%%%%%%%%%%%%%%%%%%%%%%%%%%%%%%%%%%%%%%%%%%%%%%%%%%%%%%%%%%%%%%%%%%%%%%%%%%%%%
% Theorem environments
%%%%%%%%%%%%%%%%%%%%%%%%%%%%%%%%%%%%%%%%%%%%%%%%%%%%%%%%%%%%%%%%%%%%%%%%%%%%%%%%

\newtheorem{theorem}{Theorem}
\newtheorem{lemma}{Lemma}
\newtheorem{proposition}{Proposition}
\newtheorem{corollary}{Corollary}
\newtheorem{definition}{Definition}
\newtheorem{remark}{Remark}

%%%%%%%%%%%%%%%%%%%%%%%%%%%%%%%%%%%%%%%%%%%%%%%%%%%%%%%%%%%%%%%%%%%%%%%%%%%%%%%%
% Notation
%%%%%%%%%%%%%%%%%%%%%%%%%%%%%%%%%%%%%%%%%%%%%%%%%%%%%%%%%%%%%%%%%%%%%%%%%%%%%%%%

\newcommand{\Enc}{\mathsf{Enc}}
\newcommand{\OKVS}{\text{RB-OKVS}}
\newcommand{\F}{\mathbb{F}}

\newcommand{\sender}{\mathcal{S}}
\newcommand{\receiver}{\mathcal{R}}

%%%%%%%%%%%%%%%%%%%%%%%%%%%%%%%%%%%%%%%%%%%%%%%%%%%%%%%%%%%%%%%%%%%%%%%%%%%%%%%%
% Algorithmic commands
%%%%%%%%%%%%%%%%%%%%%%%%%%%%%%%%%%%%%%%%%%%%%%%%%%%%%%%%%%%%%%%%%%%%%%%%%%%%%%%%

\algnewcommand{\algorithmicinput}{\textbf{Input:}}
\algnewcommand{\algorithmicoutput}{\textbf{Output:}}
\algnewcommand{\Input}{\item[\algorithmicinput]}
\algnewcommand{\Output}{\item[\algorithmicoutput]}

% 기존 정의의 \triangleright는 math mode가 아니므로 아래처럼 수정
\algnewcommand{\InlineComment}[1]{%
  \hspace{1em}$\triangleright$ #1%
}

%%%%%%%%%%%%%%%%%%%%%%%%%%%%%%%%%%%%%%%%%%%%%%%%%%%%%%%%%%%%%%%%%%%%%%%%%%%%%%%%
% Paragraph heading
%%%%%%%%%%%%%%%%%%%%%%%%%%%%%%%%%%%%%%%%%%%%%%%%%%%%%%%%%%%%%%%%%%%%%%%%%%%%%%%%

\makeatletter
\newcommand{\paragrap}[1]{%
  \par\addvspace{1.0ex}%
  \noindent\textbf{#1.\ }%
  \nobreak\ignorespaces
}
\makeatother

%%%%%%%%%%%%%%%%%%%%%%%%%%%%%%%%%%%%%%%%%%%%%%%%%%%%%%%%%%%%%%%%%%%%%%%%%%%%%%%%
% Draft comments
%%%%%%%%%%%%%%%%%%%%%%%%%%%%%%%%%%%%%%%%%%%%%%%%%%%%%%%%%%%%%%%%%%%%%%%%%%%%%%%%

\def\draft{1}

\ifnum\draft=1
  \newcommand{\kwag}[1]{%
    \textcolor{blue}{$\langle\!\langle$ Kwag: `#1'' $\rangle\!\rangle$}%
  }
  \newcommand{\son}[1]{%
    \textcolor{magenta}{$\langle\!\langle$ Son: `#1'' $\rangle\!\rangle$}%
  }
  \newcommand{\kwon}[1]{%
    \textcolor{purple}{$\langle\!\langle$ Kwon: ``#1'' $\rangle\!\rangle$}%
  }
\else
  \newcommand{\kwag}[1]{}
  \newcommand{\son}[1]{}
  \newcommand{\kwon}[1]{}
\fi

%%%%%%%%%%%%%%%%%%%%%%%%%%%%%%%%%%%%%%%%%%%%%%%%%%%%%%%%%%%%%%%%%%%%%%%%%%%%%%%%
% Table highlighting and additional notation
%%%%%%%%%%%%%%%%%%%%%%%%%%%%%%%%%%%%%%%%%%%%%%%%%%%%%%%%%%%%%%%%%%%%%%%%%%%%%%%%

\newcommand{\BEST}[1]{\cellcolor{magenta!20}#1}
\newcommand{\SECOND}[1]{\cellcolor{blue!20}#1}
\newcommand{\NA}{--}

\newcommand{\func}[1]{\mathcal{F}\_{{\sf #1}}}
\newcommand{\Z}{\mathbb{Z}}
\newcommand{\R}{\mathbb{R}}
\newcommand{\Q}{\mathbb{Q}}
\renewcommand{\vec}[1]{\bm{#1}}
\newcommand{\mat}[1]{\bm{#1}}

\newcommand{\red}[1]{{\color{red}#1}}
\newcommand{\blue}[1]{{\color{blue}#1}}
\newcommand{\magenta}[1]{{\color{magenta}#1}}

\newcommand{\twolines}[2]{\makecell[c]{#1\\#2}}


%-------------------------------------------------------------------------------
\begin{document}
%-------------------------------------------------------------------------------

%don't want date printed
\date{}

% make title bold and 14 pt font (Latex default is non-bold, 16 pt)
% \title{\Large \bf {\red{Efficient Private Set Operations \\ from Batch-Homomorphic OKVS Decoding}}}
\title{\Large \bf {\red{Efficient Secret-Shared Membership Tests for Private Set Operations \\ via Batched Homomorphic OKVS Decoding}}}

%for single author (just remove % characters)
% \author{
% {\rm Your N.\ Here}\\
% Your Institution
% \and
% {\rm Second Name}\\
% Second Institution
% % copy the following lines to add more authors
% % \and
% % {\rm Name}\\
% %Name Institution
% } % end author

\maketitle

%-------------------------------------------------------------------------------
\begin{abstract}
%-------------------------------------------------------------------------------
{\red{
% Private set operations (PSOs) let two parties compute set-theoretic functionalities on private inputs while revealing nothing beyond the prescribed output. While private set intersection (PSI) has become highly efficient, many other PSOs remain significantly more expensive. The most effective general framework for such tasks is based on reverse private membership test (RPMT), but even state-of-the-art RPMT constructions rely on heavy elliptic-curve-based primitives. 

% In this work, we propose a substantially faster RPMT protocol by replacing the elliptic-curve core with RLWE-based one. 
% Our starting point is the Oblivious Key-Value Store (OKVS) based RPMT framework, whose direct
% adaptation to RLWE is obstructed by the batching structure of RLWE encryption. To address this, we introduce a batching-friendly variant of OKVS together with a homomorphic batched decoding procedure. 
% We believe that this batching-friendly OKVS and its homomorphic decoding process may be of independent interest.
Private set operations (PSOs) allow two parties to compute set-theoretic functionalities on private inputs without revealing information beyond the prescribed output. A common building block for PSOs is reverse private membership test (RPMT), but RPMT reveals the membership vector to the receiver and thus at least the intersection cardinality during execution. Secret-shared private membership test (ssPMT) instead outputs XOR shares of the membership vector, enabling subsequent OT or secure-computation modules to realize various PSOs without revealing individual membership bits or the intersection cardinality. However, existing efficient ssPMT constructions still incur substantial per-execution communication and computation costs.

In this work, we propose a substantially faster ssPMT protocol by replacing the elliptic-curve core with RLWE-based one. 
Our starting point is the Oblivious Key-Value Store (OKVS) based RPMT framework, whose direct adaptation to RLWE is obstructed by the batching structure of RLWE encryption. To address this, we introduce a batching-friendly variant of OKVS together with a homomorphic batched decoding procedure. 
We believe that this batching-friendly OKVS and its homomorphic decoding process may be of independent interest.
\kwag{성능 수치 확정되면 성능 관련 문단 작성}
% For a set size $2^{20}$, our RPMT-based PSO protocols take only about $3$ seconds over LAN network and $120$-$138$MB communication, whose running time is comparable to state-of-the-art PSI. Compared to state-of-the-art PSI-Cardinality and PSI-Card-SUM, this is up to \(13.0\times\) speedups. Compared to state-of-the-art PSU, this is up to \(3.0\times\) smaller communication while achieving comparable computational cost, which results in up to $3.9\times$ faster running time over WAN. 
}}
\end{abstract}


%-------------------------------------------------------------------------------
%%%%%%%%%%%%%%%%%%%Intro%%%%%%%%%%%%%%%%%%%%%%%
\section{Introduction}
% Private set operations (PSOs) allow two parties to compute set-theoretic functionalities on their private inputs without revealing additional information beyond the prescribed output. 
% 이러한 PSO 프로토콜의 발전에 큰 영향을 미친 것은 RPMT framework~\cite{KRTW19}로, 집합 \(X\)를 가진 receiver는 sender의 각 원소 \(y\)에 대해 membership bit \(\vec{1}[y\in X]\)를 얻고, sender는 아무것도 얻는 것이 없는 구조이다. RPMT construction은 membership vector를 OT 또는 secure computation 모듈과 결합하여 다양한 PSO 프로토콜로 설계 가능한 장점이 있으며 꾸준히 발전되었다~\cite{ZCL+23,CZZ+24}.
Private set operations (PSOs) allow two parties to compute set-theoretic functionalities on their private inputs without revealing additional information beyond the prescribed output. One framework that has significantly shaped the development of PSO protocols is the reverse private set membership test (RPMT) framework~\cite{KRTW19}, in which a receiver holding a set \(X\) obtains the membership bit \(\vec{1}(y\in X)\) for each element \(y\) held by the sender, while the sender obtains nothing. RPMT constructions have been continuously developed because the resulting membership vector can be combined with OT or secure-computation modules to construct various PSO protocols~\cite{ZCL+23,CZZ+24}.

% 그러나 RPMT 기반 PSO 프로토콜에서는 PSO가 종료되기 전에 receiver가 중간 membership indication vector를 획득한다. 이에 따라 receiver는 정해진 output을 받기 전에 최소한 intersection cardinality를 알 수 있으며, 이러한 중간 정보 노출은 \cite{JSZG24}에서 \emph{during-execution leakage} (DEL)로 지적되었다.\footnote{\cite{PGT26}은 semi-honest setting에서 enhanced PSU와 standard PSU가 동등하다고 주장하며 DEL의 의미에 의문을 제기하였으나, \cite[Appendix~B]{JSZG24}는 해당 동등성 주장에 재반박하였다.} .
% PSU들의 경우 이러한 DEL 공격을 방어하기 위한 enhanced PSU의 concept 및 protocol들이 제안되고 있으나 ~\cite{JSZG24,TBZ+25}, DEL-leaky한 방향에 비해 performance가 꽤 안좋은 편임. 최근 DEL을 허용하는 상황에서는 상당히 빠른 work~\cite{PG26}도 제안되어 이러한 간극이 더 큰 상황임.
However, in RPMT-based PSO protocols, the receiver obtains an intermediate membership indication vector before the PSO protocol terminates. Consequently, the receiver can learn at least the intersection cardinality before obtaining the prescribed output, and such intermediate information leakage was identified as \emph{during-execution leakage} (DEL) in \cite{JSZG24}.\footnote{\cite{PGT26} argued that enhanced PSU and standard PSU are equivalent in the semi-honest setting, thereby questioning the significance of DEL, whereas \cite[Appendix~B]{JSZG24} subsequently rebutted this equivalence claim.} For PSU, enhanced PSU concepts and protocols have been proposed to prevent such DEL attacks~\cite{JSZG24,TBZ+25}; however, their performance is considerably worse than that of DEL-leaky approaches. More recently, a substantially faster PSU has been proposed in a setting that allows DEL~\cite{PT26}, further widening this performance gap.

% 한편, PSU 뿐만 아니라 PSO 전반에 대해 DEL에 안전한 framework에 대한 질문이 자연스럽게 떠오르고, secret-shared membership test (ssPMT)~\cite{WCS+26}가 주목받았다\footnote{사실 \cite{WCS+26}은 more generality를 위해 permutation을 추가하여 pbssPMT라고 이름붙였으나, 범용 PSO에는 ssPMT가 더 이해할 수 있음.}:
% ssPMT는 membership vector \(\vec b\)를 한 당사자에게 직접 공개하는 대신, 다음 조건을 만족하는 XOR share \(\vec b_S\)와 \(\vec b_R\)를 출력한다 s.t $\vec b_S\oplus\vec b_R=\vec b$; note that 어느 당사자도 개별 membership bit를 알 수 없어 DEL에 취약하지 않다. 이러한 security aspect의 장점 외에도, RPMT 기반의 경우 정의 자체로부터 intersection cardinality를 항상 receiver에게 알려주기 때문에 extended operation에도 이러한 제약이 걸리지만, ss-PMT는 어느 당사자에게도 먼저 공개하지 않고 후속 PSO 또는 generic secure-computation 모듈로 전달할 수 있다.
Meanwhile, beyond PSU, the question of constructing a DEL-secure framework for PSOs in general naturally arises, and the secret-shared membership test (ssPMT)~\cite{WCS+26} has attracted attention.\footnote{In fact, \cite{WCS+26} proposed more generalized concept dubbed {\em permuted-batch} ssPMT (pbssPMT) to support PSU at one stroke, but we focus on its secret-shared core.} Instead of directly revealing the membership bit $\vec 1(y\in X)$ to one party, ssPMT outputs XOR shares \(b_S\) and \(b_R\) such that $b_S \oplus b_R = \vec 1(y\in X).$
Thus, neither party learns any individual membership bit, and the protocol is not vulnerable to DEL. 
In addition to this security advantage, ssPMT has functional advantage also than RPMT that inevitably reveals the intersection cardinality to the receiver: The ssPMT can pass the membership results to a subsequent PSO or generic secure-computation module without revealing them to either party.

% \cite{WCS+26}이 제안한 ss-PMT는 상당히 efficient한 online phase를 달성하였으나, 상당히 무거운 offline phase를 필요로 함. 비용이 큰 연산을 input-independent offline phase로 이동할 수 있다는 점은 online latency를 줄이는 데 유용하지만, 해당 offline phase는 한 번의 setup 이후 여러 실행에서 재사용되는 것이 아니다. Offline에서 생성된 correlated state는 각 PSO 실행에서 한 번씩 소비되므로, 새로운 실행마다 다시 생성하고 통신해야 한다. 따라서 실행당 필요한 전체 communication 및 computation cost는 여전히 크며, practical한 관점에서 개선의 여지가 상당하다.
The ssPMT proposed by \cite{WCS+26} achieves a highly efficient online phase, but requires a substantial offline phase. Moving expensive operations to an input-independent offline phase is useful for reducing online latency,  but this offline phase is not a one-time setup that can be reused across multiple executions. The correlated state generated offline is consumed once in each PSO execution and must therefore be regenerated and communicated for every new execution. Consequently, the total communication and computation costs required per execution remain large, leaving substantial room for improvement from a practical perspective.


%%%%%%%%%%%%%%%%%%%%%%%%%%%%%%%%%%%%%%%%%%%%%%%%%%%%%%
\subsection{Our Contributions}
Our goal is to replace the expensive elliptic-curve core in prior oblivious key-value store (OKVS) based ssPMT with a much lighter ring learning with errors (RLWE) based one. 
%\blue{This change not only enables efficient batching but also gives a post-quantum-secure RPMT core under the standard RLWE assumption.} 
Our starting point is the RPMT framework of Zhang et al.~\cite{ZCL+23}, whose core operation is OKVS decoding over ElGamal-encrypted state. 
The main challenge for adapting their framework with RLWE-based encryption is the batching property of RLWE-based encryption: in RLWE-based encryption, multiple messages are encrypted into a single ciphertext, and this property makes RLWE-based encryption incompatible to the prior idea to build RPMT.

\paragrap{Batching-friendly homomorphic OKVS decoding}
We present \emph{RSB-OKVS}, a batching-friendly variant of Random-Band OKVS~\cite{BPSY23} tailored to the SIMD layout of RLWE encryption. 
% Concretely, we propose a matrix construction obtained by a suitable linear transformation of the original RB matrix, so that the encoding procedure and failure-probability analysis of RB-OKVS can be carried over, while restructuring the decoding pattern so that it aligns with RLWE batching.
Concretely, we construct RSB-OKVS by applying a column permutation, equivalently a linear transformation, to the original RB matrix. This preserves the encoding procedure and failure-probability analysis of RB-OKVS, while restructuring the decoding pattern so that it aligns with RLWE batching.
By our efficient \emph{sequencing} technique and  \emph{homomorphic batched decoding} procedure for RSB-OKVS, we resolve the above problems and allow OKVS decoding to be evaluated efficiently using lightweight slot-wise homomorphic operations without ciphertext rotations. 

\paragrap{Implementation and Performance}
We implement our ssPMT-based PSOs and evaluate them for input sizes $2^{16}$--$2^{22}$ under bandwidths from $10$Gbps to $10$Mbps, comparing with prior protocols.
For a direct comparison of ssPMT protocol, we focus on \cite{WCS+26}, which produces secret-shared membership results without revealing the intermediate membership vector. At \(n=2^{20}\), our ssPMT reduces the total communication from \(3731.01\) MB to \(199.17\) MB, an \(18.73\times\) reduction. It also reduces the total runtime by \(2.96\times\), \(2.86\times\), and \(8.88\times\) over LAN, \(1\) Gbps, and \(100\) Mbps, respectively. These results show that our construction substantially lowers the per-execution cost of secret-shared membership testing, while retaining competitive online latency.

% For PSU, although the runtime of our protocol is about $1.5\times$ slower than the state-of-the-art PSU protocol~\cite{PT26} on LAN, it reduces communication by up to $3.4\times$ and thus performs better on slower networks; on sub-100Mbps networks, it achieves up to $3.06\times$ speedup over~\cite{PT26}. 
% For PSI-Card and PSI-Card-Sum, our protocol achieves the best runtime for almost every input sizes in our tested network environments. Compared with the state-of-the-art protocols~\cite{CZZ+24}, our protocols achieves up to $13.4\times$ speed-up over LAN network, with only $1.1-3.1\times$ overheads on the communication. This results in better runtime for all 1Gbps and 100Mbps settings, and large-input 10Mbps settings.

%%%%%%%%%%%%%%%%%%%%%%%%%%%%%%%%%%%%%%%%%%%%%
\section{\blue {Technical Overview}}\label{sec:tech_overview}
This section provides a technical overview of main part in our new construction of secret-shared Private Membership Test (ssPMT) whose simplified definition is given as \Cref{fig:F-sspmt-informal}, using OKVS and RLWE-batched PKE.

\begin{figure}[htbp]
\centering
\fbox{%
\begin{minipage}{0.95\linewidth}
\textbf{Parameters:} Sender $\mathcal{S}$, Receiver $\mathcal{R}$, set sizes $n_y$ and $n_x$.\\
\textbf{Input:} $X\subset\{0,1\}^*$ of $\receiver$, and $Y = \{y_i\} \subset\{0,1\}^*$ of $\sender$\\
\textbf{Functionality:} Define $\vec b = (b_i = \vec 1(y_i \in X))_{i \in [n_y]}$, and return $\vec b_\sender$ and $\vec b_\receiver$ to each party such that $\vec b_\sender \oplus \vec b_\receiver = \vec b$.
\end{minipage}
}
\caption{(Simplified) Ideal functionality $\mathcal{F}_{ssPMT}$.}
\label{fig:F-sspmt-informal}
\end{figure}

The readers are assumed to be familiar with the batching nature (or SIMD packing) and additive homomorphic property of RLWE-PKE, and the random-band matrix based OKVS construction~\cite{BPSY23}, whose review is given in \Cref{sec:prelim}.

% \begin{figure}[htbp]
% \centering
% \fbox{%
% \begin{minipage}{0.95\linewidth}
% \textbf{Parameters:} Sender $\mathcal{S}$, Receiver $\mathcal{R}$, set sizes $n_y$ and $n_x$.\\
% \textbf{Input:} $X\subset\{0,1\}^*$ of $\receiver$, and $Y = \{y_i\} \subset\{0,1\}^*$ of $\sender$\\
% \textbf{Functionality:} Return $\vec b = (b_i = \vec 1(y_i \in X))_{i \in [n_y]}$ to $\mathcal{R}$.
% \end{minipage}
% }
% \caption{Ideal functionality $\mathcal{F}_{RPMT}$.}
% \label{fig:F-mq-rpmt}
% \end{figure}

\paragrap{Review of ZCLZL encrypted PMT}
We briefly review the starting point of our work, ZCLZL~\cite{ZCL+23}. It was originally aimed to construct private set union, but it suffices to review its technical core that tests membership in an encrypted state using OKVS and ElGamal PKE.\footnote{\cite{ZCL+23} also proposed SKE based instantiation, but we only focus on PKE version because SKE version has poor performance, due to expensive preprocessing to generate Beaver Triple.}

The main idea is simple: the receiver $\receiver$ creates $n_x = |X|$ different ciphertexts $c_i$ of some indicator string ${\sf IS}$, and OKVS-encodes them by using each $x_i$ as a key and the corresponding ciphertext $c_i$ as its value. The sender $\sender$ then receive the OKVS-encoding, and decodes it at each $y_j \in Y$ to obtain $c_j^*$. Then, when $y_j \in X$, $c_j^*$ is a ciphertext of ${\sf IS}$; otherwise, it is a ciphertext that contains another value than ${\sf IS}$. 
% Although original \cite{ZCL+23}은 $c_j^*$를 곧바로 receiver에게 돌려보내주고\footnote{사실 엄밀히는 re-randomize가 필요한데, 우리 construction도 이건 마찬가지라 main body에서 설명한다.} 바로 decrpyt해서 receiver가 membership bit $\vec b$를 알아내는 so-called Reverse-PMT construction을 제안함. 이는 \cite{JSZG24,TBZ+25}에서 지적받은 대로 DEL에 취약한 버전이지만, 위 아이디어를 조금만 tweak해서 ss-PMT로 adapt할 수 있다, using ElGamal PKE의 additive homomorphic 성질: $\sender$가 some uniform mask $r_j$로 $c_j^*$를 masking해서 보내면, $\receiver$의 $c_j^*$ decryption 결과 $d_j$와 $r_j$는 when $y_j \in X$, $d_j - r_j = \mathsf{IS}$, otherwise $d_j - r_j \neq \mathsf{IS}$을 만족함.
% 이후 Two parties는 generic two-party computation을 이용하여 $\vec 1(r_j = d_j - \mathsf{IS})$를 secret-shared 상태로 계산하여, formalized by $\mathcal{F}_{\mathsf{ss-peqt}}$~\cite{TBZ+25}, ss-PMT 결과물을 얻을 수 있다.
% 이 내용을 \Cref{fig:zcl_pke_framework}로 정리했고, 우리는 이것을 starting point로 삼는다.
The original construction of \cite{ZCL+23} returns $c_j^*$ to the receiver,\footnote{Strictly speaking, $c_j^*$ must
first be re-randomized. The same requirement arises in our construction and is discussed in the main body.} who decrypts it and directly learns the membership vector $\vec b$, thereby realizing the so-called Reverse-PMT functionality. As pointed out in \cite{JSZG24,TBZ+25}, this construction is vulnerable to DEL. 
Nevertheless, the same basic idea can be adapted to ssPMT with a minor modification by exploiting the additive homomorphic property of ElGamal PKE.
Specifically, the sender homomorphically masks each $c_j^*$ with a uniform value $r_j$ before returning it. 
Let $d_j$ denote the receiver's decryption result. Then, $ d_j-r_j=\mathsf{IS} \quad\text{if }y_j\in X,$ whereas $d_j-r_j\neq\mathsf{IS} \quad\text{if }y_j\notin X.$
The parties then use a generic two-party computation to evaluate $\mathbf{1}[r_j=d_j-\mathsf{IS}]$ in secret-shared form, as formalized by $\mathcal{F}_{\mathsf{ss\mbox{-}peqt}}$~\cite{TBZ+25}, and thereby obtain the ssPMT output. We summarize this adaptation in \Cref{fig:zcl_pke_framework} and use it as the starting point of our construction.



\begin{figure}[htbp]
\centering
\fbox{%
\begin{minipage}{0.95\linewidth}
\textbf{Input :} The receiver $\receiver$ with $X$, and the sender $\sender$ with $Y=\{y_j\}_{j\in[n_y]}$
% \textbf{Output :} $\receiver$ obtains $\vec b = \{b_i = \vec 1(y_i \in X)\}_{i \in [n_y]}$
\begin{enumerate}
\item $\mathcal{R}$ samples a random indicator ${\sf IS}$, and  encrypts it $n_x$ times to have ciphertexts $\{c_i\}_{i\in[n_x]}$. 
\item $\mathcal{R}$ computes the OKVS encoding $\vec p$ of the key-value table $\{(x_i, c_i)\}_{i\in[n_x]}$, and send $\vec p$ to $\mathcal{S}$.
\item $\mathcal{S}$ decodes the OKVS encoding $\vec p$ for each $y_j$ to have a ciphertext $c_j^*$, re-randomizes $c_j^*$ by adding an encryption of zero.
\item $\mathcal{S}$ masks each $c_j^*$ with a random masking $r_j$, and send back it to $\mathcal{R}$, who decrypts it to have $d_j$.
\item Two parties engage (secret-shared) private equality test on input $r_j$ and $d_j - \mathsf{IS}$, to have $b_{\sender, j}$ and $b_{\receiver, j}$ s.t. $b_{\sender, j} \oplus b_{\receiver, j} = \vec 1(r_j = d_j - \mathsf{IS})$.
\end{enumerate}
\end{minipage}
}
\caption{An adaptation of \cite{ZCL+23} for ssPMT, based on OKVS and a re-randomizable PKE.} 
\label{fig:zcl_pke_framework}
\end{figure}

% \subsection{From PKE-Based RPMT of ZCLZL \\ to Homomorphic OKVS Decoding}
\subsection{RLWE-PKE instantiation of ZCLZL}
Our goal is to replace the ElGamal PKE in the ssPMT variant of \cite{ZCL+23} (\Cref{fig:zcl_pke_framework}) by the RLWE-based PKE. As a starting obstacle, note that the original construction implicitly assumes that one message corresponds to one ciphertext, which makes it possible to set the key-value table as $\{(x_i, c_i)\}$. However, RLWE-PKE instantiation immediately confronts a difficulty; it cannot associate one ciphertext to each key $x_i$, because one RLWE ciphertext contains multiple messages. 
 
\paragrap{Switching OKVS encoding and encryption}
We resolve this issue by switching the order of OKVS encoding and encryption. Namely, we first OKVS-encode the key-value table $\{(x_i, \mathsf{IS})\}$, where all $x_i$ are assigned the same value ${\sf IS}$, to obtain $\vec{p}$, and then encrypt the vector $\vec p$. The ZCLZL construction is still valid despite of this switching, since RLWE-PKE is well known to have an additive homomorphic property. More precisely, one can homomorphically perform OKVS decoding on $y$, equivalently, the inner-product computation $\langle \mathsf{row}(y), \vec{p} \rangle$ between some sparse vector $\mathsf{row}(y)$ and $\vec{p}$. Here, $\mathsf{row}(y)\in\mathbb{F}^{m}$ denotes the sparse OKVS decoding row associated with $y$, satisfying $\mathsf{Decode}(y,\vec{p})=\langle\mathsf{row}(y),\vec{p}\rangle$. Therefore, one obtains the OKVS decoding result in the encrypted state.

In fact, this idea of switching the order of encoding and encryption had already been considered in \cite{ZCL+24}, but with ElGamal PKE. In that case, the homomorphic evaluation of $\langle \mathsf{row}(y), \vec{p} \rangle$ is very trivial. More specifically, writing $\mathsf{row}(y) = (r_i)$, we have
\begin{equation}\label{eq:decodeinnerprod}
\langle \mathsf{row}(y), \vec{p} \rangle = \sum_{r_i \neq 0} r_i \cdot p_i,
\end{equation}
and the sender can simply take the corresponding ElGamal ciphertext of each $p_i$ and multiply it by the coefficient $r_i$.

However, when encryption is performed with RLWE-PKE, $\vec{p}$ is encrypted in a batched sense. More precisely, writing $\vec{p} = (\vec{p}_1 \,\|\, \cdots \,\|\, \vec{p}_b)$ for $b = \lceil m / N \rceil$, each length-$N$ vector $\vec{p}_i$ is encrypted into a single ciphertext $c_i$. This is incompatible with the entry-wise computation used in ElGamal PKE.


\paragrap{OKVS decoding as a matrix--vector product}
Therefore, the main technical problem is to compute the encrypted inner product~\eqref{eq:decodeinnerprod} where $\vec{p}$ is encrypted in a batched manner. For this, note that the sender after all needs to decode the OKVS \emph{for all} $y \in Y$. In this view, OKVS decoding actually corresponds to a matrix--vector product $\mat{M} \cdot \vec{p}$, where the row of $\mat{M}$ corresponding to $y$ is $\mathsf{row}(y)$ and $\vec{p}$ is the OKVS encoding vector. At this point, also note that we exclusively consider the random-band (RB) matrix-based OKVS construction of~\cite{BPSY23}, in which each $\mathsf{row}(y)$ consists of a random band---namely, a short sequence of entries of fixed length---starting at a random position.

Then, considering that RLWE batching and slot-wise operations are particularly well suited for diagonal matrix--vector product, it is natural to realign the rows of $\mat{M}$ so that, with respect to their band starting positions, the matrix becomes (almost) diagonal; we can then formulate homomorphic OKVS decoding to a matrix--vector product where the matrix consists of consecutive diagonals. 
%삭제 3
% See \Cref{fig:RSB matrix} for a graphical illustration.
% %%%%%%%%%%%%%%%%%%%%%그림%%%%%%%%%%%%%%%%%%%
% \begin{figure}[h]
%     \centering
%     % 왼쪽 그림
%     \begin{subfigure}{0.23\textwidth}
%         \centering
        
%             \includegraphics[width=\textwidth]{images/bf sorting.jpg}
%         \caption{Random band matrix~\cite{BPSY23}}
%         \label{fig:bfsorting matrix}
%     \end{subfigure}
%     \hfill
%     % 오른쪽 그림
%     \begin{subfigure}{0.23\textwidth}
%         \centering
%         \raisebox{0.4cm}{%
%         \includegraphics[width=\textwidth]{images/af sorting.jpg}
%         }
%         \caption{Diagonalization}
%         \label{fig:RB-decoding-a}
%     \end{subfigure}
%     \caption{RB-OKVS decoding is a matrix--vector product with the matrix in (a). Reordering rows to an (almost) diagonal matrix in (b) makes it more compatible with RLWE slot-wise operations.}
%     \label{fig:RSB matrix}
% \end{figure}
% %%%%%%%%%%%%%%%%%%%%%%%%%%%%%%%%%%

% {\blue{
% \paragraph{Why a KVS suffices.}
% The above reformulation also clarifies that our protocol does
% not require the obliviousness property of an OKVS. The
% encoding vector $\vec p$ is transmitted only through
% IND-CPA-secure RLWE ciphertexts, and its decoding results
% remain encrypted until they are masked before decryption.
% Consequently, the security of our protocol does not rely on
% the plaintext encoding being indistinguishable from random.
% It suffices to use a linear KVS that correctly decodes the
% stored keys and supports the sparse inner-product
% representation above. For PMT correctness, we additionally
% require only that an absent key matches the distinguished
% indicator with negligible probability, namely,
% $
% \Pr[\mathsf{Decode}(y,\vec p)=\mathsf{IS}\mid y\notin X]
% \leq 2^{-\lambda}
% $.
% Although we instantiate the protocol using the random-band
% construction originally introduced as an OKVS, its stronger
% obliviousness property is not used in our security proof.}}

\subsection{Homomorphic RB Matrix-Vector Product}
The diagonalized random band is more homomorphic-friendly, but two problems remain.
\begin{itemize}
    \item[(1)] Each row consists of a contiguous random band of length $w$. A naive matrix--vector product would therefore still be expensive, as it would require $w$ rotated copies of the encoding vector $\vec{p}$.
    % \footnote{For a non-HE batch familiar reader, \Cref{sec:rlwe_and_batch} might help understanding.}
    \item[(2)] In fact, a fully diagonalized matrix is rarely achievable, since the starting position of each random band is determined randomly from the underlying items in $Y$, which will yields collisions between the starting positions with overwhelming probability. 
    
\end{itemize}
% We believe the core technical novelty of this work is in our solutions for these issues, which are described below.

\paragrap{Random spaced bands}
Let us address the first problem, while assuming that the second problem (full diagonalization) is resolved somehow. Our idea is to replace each contiguous band by $N$-spaced entries, where $N$ is the RLWE slot size. Note that if this $N$-spacing idea works well, the encrypted computation of $\mat{M}\cdot \vec{p}$ can be carried out using only slot-wise additive homomorphic operations, without any ciphertext rotations; see \Cref{fig:wRSB matrix}.
%%%%%%%%%%%%%%%%%%Figure%%%%%%%%%%%%%%%%%%%%%%%%%
\begin{figure}[h]
    \centering
    % 왼쪽 그림
    \begin{subfigure}[h]{0.23\textwidth}
        \centering
        \includegraphics[width=\textwidth]{images/RB matrix.PNG}
        \caption{(Diagonalized) Random bands}
    \end{subfigure}
    \hfill
    % 오른쪽 그림
    \begin{subfigure}[h]{0.23\textwidth}
        \centering
        \includegraphics[width=\textwidth]{images/RSB matrix.PNG}
        \caption{(Diagonalized) Random spaced bands}
    \end{subfigure}
    \caption{Diagonalized matrices. (Each small square represents a $N \times N$ sub-block.)}
    \label{fig:wRSB matrix}
\end{figure}
%%%%%%%%%%%%%%%%%%%%%%%%%%%%%%%%%%%%%%%%%%%%%%%%%

To make this $N$-spacing idea more substantial, it should be noted that all prior OKVS-related matrix proposals~\cite{PRTY20,GPR+21,RR22,BPSY23} involve a complicated argument related to encoding/decoding efficiency and failure probability analysis, which are indeed the most complicated part when proposing a new OKVS construction. Therefore, we should provide the same argument for our new spaced matrix family. 
As a simple solution, we define the \emph{random spaced band} (RSB) matrix as a fixed column-wise permutation of a random-band matrix, where the permutation is determined by a rigorous extension of the $N$-spacing idea. More precisely, given a key set $K$, we construct the corresponding random spaced band matrix as in \Cref{alg:RSB matrix generato}, where the concrete form of the permutation matrix $\mat C$ is given in the main body as \Cref{eq:N_perm_mat}.

%%%%%%%%%%or%%%%%%%%%%%%
\begin{algorithm}[htbp]
\caption{\textsf{RSB-OKVS.GenMat}}
\begin{algorithmic}[1]
\Input A key set $K = \{k_i\}_{i \in [n]}$ and a spacing length $N$
\Output Spaced random matrix $\mat{\hat{M}}$
\State Generate a random band matrix $\mat M$ using $K$
\State \Return $\mat{\hat{M}}=\mat M \cdot \mat C_N$ where $\mat C_N$ is the $N$-spacing matrix defined as \Cref{eq:N_perm_mat}.
\end{algorithmic}\label{alg:RSB matrix generato}
\end{algorithm}
%%%%%%%%%%%

Thanks to this column permutation-based definition, encoding with respect to the RSB matrix can be carried out easily by using the efficient encoding algorithm for the RB matrix~\cite{BPSY23}, as in \Cref{alg:RSBEncoding}. Furthermore, the full rank failure-probability analysis for the RB matrix in~\cite{BPSY23} remains valid for the RSB matrix as well. More precisely, if the row size $m\approx(1+\varepsilon) n$ and the band width $w$ provide failure probability $2^{-\lambda}$ for the random-band matrix, then of course the RSB matrix produced by \Cref{alg:RSB matrix generato} does as well. Therefore, the methodology used in \cite{BPSY23} to choose $(\varepsilon,w)$ (explained in \Cref{sec:OKVS experiment}) is still valid in our case, and in fact we set the parameters in this way.

%%%%%%%%%%%%%
\begin{algorithm}[htbp]
\caption{\textsf{RSB-OKVS.Encode}}
\begin{algorithmic}[1]
\Input A key-value set $\{(k_i, v_i)\}_{i \in [n]}$, a spacing length $N$
\Output Encoding $\hat{\vec p}$
\State Generate a random band matrix $\mat M$ using $K = \{k_i\}_{i \in [n]}$
\State Solve the linear system $\mat M\vec p = \vec v$ (using the efficient method of~\cite{BPSY23})
\State \Return $\hat{\vec p}=\mat C_N^{-1} \vec p$ where $\mat C_N$ is defined as \Cref{eq:N_perm_mat}.
\end{algorithmic}\label{alg:RSBEncoding}
\end{algorithm}
%%%%%%%%%%%

\begin{remark}
    % 우리는 OKVS의 correctness 성질만 필요하고, security--or \em{oblivious}-- property는 전혀 필요하지 않다, because the sender at most observes its encryption. Nonetheless, we believe any \em{oblivious} property of RB-OKVS also holds for our RSB-OKVS also, and it might be useful for any OKVS applications where such property is necessary~\cite{ZCL+23,RS21}.

    Although our protocol uses only the correctness property of the OKVS and does not rely on its obliviousness, any distributional obliviousness property that is invariant under a fixed bijective relabeling of the encoding coordinates is preserved by our RSB transformation. In particular, the random-encoding property required in \cite{ZCL+23,RS21} is preserved.
\end{remark}


As one downside of the column-permutation-based definition, a contiguous random band of length $w$ is shifted by one position each time it crosses a row boundary, as illustrated by the green super-diagonal in \Cref{fig:wRSB matrix2}. Because of this, homomorphic OKVS decoding can no longer be carried out using only the ciphertexts of $\vec p$. More precisely, at most $w-1$ such super-diagonals can arise (when the starting position is placed near the end of a row), and therefore we need at most $w-1$ additional ciphertexts of such shifted copies, or equivalently, $w-1$ homomorphic rotations. Nevertheless, the resulting matrix--vector product is still far more efficient than the original RB-matrix-based computation. More details can be found in \Cref{subsec:RSB OKVS}.
%%%%%%%%%%%%%%%%%%Figure%%%%%%%%%%%%%%%%%%%%%%%%%
\begin{figure}[h]
    \centering
    % 왼쪽 그림
    \begin{subfigure}[h]{0.23\textwidth}
        \centering
        \includegraphics[width=\textwidth]{images/RB matrix2.jpg}
        \caption{(Diagonalized) Random bands}
    \end{subfigure}
    \hfill
    % 오른쪽 그림
    \begin{subfigure}[h]{0.23\textwidth}
        \centering
        \includegraphics[width=\textwidth]{images/RSB matrix2.jpg}
        \caption{(Diagonalized) Random spaced bands}
    \end{subfigure}
    \caption{Under our RSB matrix definition based on the column permutation, each row boundary cross of the $N$-spacing band induces one right shift. This is represented as a green super-diagonal in (b).}
    \label{fig:wRSB matrix2}
\end{figure}
%%%%%%%%%%%%%%%%%%%%%%%%%%%%%%%%%%%%%%%%%%%%%%%%%
%%%%%%%%%%%%%%%%%%Figure%%%%%%%%%%%%%%%%%%%%%%%%%
\newlength{\seqfigheight}
\setlength{\seqfigheight}{3.9cm} 

\begin{figure*}[t]
    \centering
    \begin{subfigure}[t]{0.26\textwidth}
        \centering
        \includegraphics[width=\linewidth,height=\seqfigheight,keepaspectratio]{images/seq1.png}
        \caption{Original Random Spaced Matrix}
        \label{fig:sequencing_overview_a}
    \end{subfigure}\hfill
    \begin{subfigure}[t]{0.465\textwidth}
        \centering
        \includegraphics[width=\linewidth,height=\seqfigheight,keepaspectratio]{images/seq2.png}
        \caption{Diagonal $m\times m$ matrices}
        \label{fig:sequencing_overview_b}
    \end{subfigure}\hfill
    \begin{subfigure}[t]{0.25\textwidth}
        \centering
        \includegraphics[width=\linewidth,height=\seqfigheight,keepaspectratio]{images/seq3.png}
        \caption{Block-diagonal $N\times m$ matrices}
        \label{fig:sequencing_overview_c}
    \end{subfigure}

    \caption{A toy-example of our matrix sequencing, with $n = 12, m = 15, w = 3$ and $N = 4$. (b) corresponds to the naive diagonalization sequencing, and (c) is an example of our proposed sequencing. Since each nonempty $N\times N$ block translates into one slot-wise multiplication, (b) needs 27 such multiplications, whereas (c) needs only 16.}
    \label{fig:sequencing_overview}
\end{figure*}
%%%%%%%%%%%%%%%%%%%%%%%%%%%%%%%%%%%%%%%%%%

\paragrap{Matrix sequencing to handle collisions}
It remains to address the second issue, namely collisions between band starting positions. The most naive fix would be to build multiple $m \times m$ diagonalized matrices and distribute colliding rows across them, as \Cref{fig:sequencing_overview_b}. However, it can be easily argued that the number of $m \times m$ matrices obtained in this way is $O(\log n)$~\cite{RS98}, which results in up to $O(\log n)$ multiplications of RSB matrix and OKVS encoding vector.\footnote{For $m = O(n)$, the maximum number of rows that share the same starting position is $O(\log n)$ with high probability by \cite{RS98}.}

We actually propose a more efficient method by utilizing the RSB matrix structure. Observe that if two rows have the same starting position {\em modulo $N$}, then their bands lie on the same diagonal within an $N \times N$ submatrix after diagonalization. Motivated by this observation, we design a procedure that packs the rows according to their {\em modulo $N$} starting positions to construct multiple $N \times m$ matrices in a diagonal-block shape; see \Cref{fig:sequencing_overview_c} for an example. 
We achieve only constant times of computational overhead, compared to the ideally diagonalizable case that involves only one $m\times m$ matrix--vector product. More details can be found in ~\Cref{subsec:sequencing}.

% \subsection{Handling Sequencing Layout}\son{제목 고민?}
\subsection{Padding Sequencing Outputs}
The random-spaced-band transformation and matrix sequencing described above enable an efficient homomorphic evaluation of $\mat M\vec p$. In other words, they provide an RLWE-based realization of the encrypted OKVS decoding step 3 in \Cref{fig:zcl_pke_framework}. However, sequencing introduces a new security issue absent from the entry-wise ElGamal instantiation, which must be resolved to complete the ssPMT construction.

The query matrix $\mat M$, whose rows are generated from the sender's input $Y$, is partitioned by the sequencing
algorithm into $L$ layer matrices of dimension $N\times m$.
Homomorphic decoding consequently produces one RLWE ciphertext for each layer, for a total of $L$ ciphertexts.
Since the collision pattern of the query-row starting positions depends on $Y$, the resulting layer count $L$ may
also depend on the sender's input. Therefore, returning exactly these $L$ ciphertexts would reveal a function of $Y$
through the response length. 

To hide the input-dependent layer count, we set a public upper bound $L'=L'(n_y,N,m,\zeta)$ of $L$, 
% If $L>L'$, the protocol returns $\bot$; otherwise, 
and let the sender pads the resulting ciphertexts up to $L'$ ciphertexts.{\footnote{This follows the same high-level rationale as cuckoo-hashing-based protocols that choose a fixed public capacity so that the hashing overflow probability is below $2^{-40}$.}}. 
This clearly incurs larger communication for returning ciphertexts and more works for ss-peqt, but does not affect the main homomorphic decoding cost; it depends only on the actual number of layers $L$. 
The concrete choice of $L'$ is analyzed in \Cref{lemma:Lprime}.

%%%%%%%%%%%%%%%%%%%%%%Pre%%%%%%%%%%%%%%%%%%%%%%%%%%
\section{Preliminaries}\label{sec:prelim}
For $n\in \mathbb{N}$, we denote the set $\{1,2,\dots,n\}$ by $[n]$.
We use $\kappa,\lambda$ to indicate the computational and statistical security parameters, respectively.
Vectors are written in lowercase bold letters (e.g., $\vec{v}$) or by its elements wrapped in parenthesis with the range of index (e.g., $(v_i)_{i\in [n]}$). Matrices are written in uppercase bold letters (e.g., $\mat{A}$), their elements are denoted by subscripts $\mat{A}_{i,j}$, or sometimes $\mat A[i][j]$. Each $i$-th row of matrix $\mat A$ is denoted by $\mat A[i]$. 
We use the indicator notation $\mathbf{1}(S)$, which equals $1$ if a statement $S$ is true.

\subsection{RLWE-based Encryptions and Batching}\label{sec:rlwe_and_batch}
We use an encryption scheme whose security is based on the hardness of Ring Learning with Errors (RLWE)~\cite{LPR10}. Its message space is $R_p := \mathbb{Z}_p[X]/(X^{N}+1)$ for some message modulus $p$ and dimension $N$, and ciphertext space is $R_q^2$ for some ciphertext modulus $q$.
For a power-of-two $N$ and $p \equiv 1~(\bmod~2N)$, the plaintext space is isomorphic to $\mathbb{Z}_p^{N}$ (as rings), and then one RLWE ciphertext can encrypt multiple integers. We call this the {\em batching} (or SIMD) property, and write the encryption by $\mathsf{RLWE.Enc}: \vec{v} \in \mathbb{Z}_p^N \mapsto c \in R_q^2.$

\paragrap{Slot-wise homomorphic operations} 
RLWE encryption is well-known to be {\em additively homomorphic}. Under the batching view, one can perform slot-wise arithmetic in the encrypted state: for ciphertexts $c_1,c_2$ encrypting $\vec{v}_1,\vec{v}_2 \in \mathbb{Z}_p^{N}$, and any plaintext vector $\vec{v}\in \mathbb{Z}_p^{N}$, we have
\begin{align*}
\mathsf{RLWE.Dec}\left(\mathsf{RLWE.CtxtAdd}(c_1,c_2)\right) =\vec{v}_1 + \vec{v}_2 \in \mathbb{Z}_p^{N} \\
  \mathsf{RLWE.Dec}\left(\mathsf{RLWE.PlainAdd}(c_1,\vec{v})\right) =\vec{v}_1 + \vec{v} \in \mathbb{Z}_p^{N}, 
\\  \mathsf{RLWE.Dec}\left(\mathsf{RLWE.PlainMult}(c_1, \vec{v})\right) =\vec{v}_1 \odot \vec{v} \in \mathbb{Z}_p^{N},
\end{align*}
where $\odot$ denotes component-wise multiplication.

\begin{remark}
RLWE encryption also supports more advanced homomorphic operations, such as ciphertext multiplication and slot rotations, but these are much more expensive than the additive operations used here. Our protocol relies only on additive homomorphic operations, which is a fundamental reason for its good performance.
\end{remark}

\paragrap{(Super-)Diagonal matrix-vector product} 
The slot-wise multiplication can be naturally understood as a diagonal matrix--vector product. More precisely, the multiplication of a diagonal matrix and a vector corresponds to a slot-wise multiplication of the diagonal and the vector. More precisely, it holds that for a diagonal matrix $\mat{D} = \mathsf{diag}(\vec d)$,
\[
  \mat{D}\cdot \vec p = \vec d \odot \vec p.
\]


For a generalization, consider a $k$-th super-diagonal matrix that has only nonzero entries $d_i$ on the position $(i, [i+k]_N)$, denoted by $\mathsf{sdiag}_k(\vec d)$. Then it holds that
\[
  \mathsf{sdiag}_k(\vec d) \cdot \vec p = \vec d \odot \rho_k(\vec p)
\]
where $\rho_k(\cdot)$ is the left rotation by $i$ slots, i.e.,
\[
  \rho_k: (p_1, \dots, p_N) \mapsto (p_{1+k}, p_{2+k}, \cdots, p_{N+k})
\] 
where all indices are taken modulo $N$ into $\{1, 2, \cdots, N\}$.



%%%%%%%%%%%%%%%%%%%%%%%%%%%%%%%%%%%%%%%%%%%%%%%%%%%%%%
% {\blue{
% \subsection{Key-Value Stores and Random-Band Constructions} \label{sec:RB-OKVS} 
% A key-value store (KVS)~\cite{GPR+21} is a data structure that compactly represents a prescribed mapping from a set of keys to their associated values. 
% \begin{definition}[Key-Value Store] A KVS is parameterized by a key universe $\mathcal K$, a finite field $\F$, and an expansion ratio $\varepsilon>0$, and consists of the following two algorithms: 
%     \begin{itemize}
%         \item $\mathsf{Encode}(T)$:  The encode algorithm receives a set of $n$ distinct key-value table 
%         $T = \{ (k_1, v_1), \dots, (k_n, v_n) \} \in (\mathcal{K} \times \F)^n$ and outputs the encoding $\vec p \in \F^m$ where $m \approx (1 + \varepsilon)\cdot n$.    
%         \item $\mathsf{Decode}(k, \vec p)$:  
%         The decode algorithm receives the encoding $\vec p \in \F^m$, a key $k \in \mathcal{K}$ and outputs some value $v \in \F$.    
%     \end{itemize}
% \end{definition}

% A KVS is correct if, for every table $T=\{(k_i,v_i)\}_{i\in[n]}$, the encoding $\vec p\leftarrow\mathsf{Encode}(T)$ satisfies $\mathsf{Decode}(k_i,\vec p)=v_i$ for all $i\in[n]$ with overwhelming probability. 

% \paragrap{Relation to oblivious key-value stores}
% It is known that if the distribution of the encoding vector is statistically indistinguishable from the uniform distribution over $\mathcal V^m$, then the corresponding OKVS is \emph{doubly oblivious}.  Our construction does not require this property.  In particular, the sender never learns the encoding vector $\vec p$ in the clear, but receives only RLWE encryptions of its packed blocks and evaluates the decoding operations homomorphically.
% Therefore, we do not require the encoding output itself to be indistinguishable from uniform.

% Similarly, some OKVS constructions require the decoding result at an absent key to be statistically uniform, a property commonly referred to as \emph{random decoding}. Our privacy argument does not rely on this full property, since the sender never observes the plaintext decoding results. 
% However, correctness still requires that an absent key decode to the distinguished indicator $\mathsf{IS}$ only with negligible probability. 우리는 $\Pr[\mathsf{Decode} (k, \vec p)= \sf{IS} | k \notin \mathcal K]\leq 2^{-\lambda}$가 되도록하는 $\ell_{\mathsf{IS}}$를 고려하여 파라미터를 잡는다.

% Accordingly, throughout this paper, we use the notation \emph{(O)KVS} as an umbrella term when a statement applies to both KVS and OKVS constructions.  More specifically, our protocol relies on the correctness of the KVS interface and the linear decoding structure of the RB-OKVS construction described below, but does not rely on the obliviousness of its public encoding.

% % Prior work commonly augments the above KVS interface with an obliviousness requirement on the public encoding.  In particular, under the standard terminology, if the encoding vector is statistically indistinguishable from the uniform distribution over $\mathcal V^m$, then the resulting construction is \emph{doubly oblivious}.

% % Our protocol does not rely on this distributional property.  The receiver never reveals the encoding vector $\mathbf p$ in the clear; instead, it sends only RLWE encryptions of its packed blocks, and the
% % sender evaluates all decoding operations homomorphically.  Therefore, the security of our protocol does not require the encoding distribution itself to be indistinguishable from uniform.

% % To distinguish the linear key-value-store interface used in our protocol from the additional obliviousness properties required in some applications, we use the term \emph{(oblivious) key-value store}, abbreviated as \emph{(O)KVS}, when referring collectively to KVS and OKVS constructions.  Our protocol specifically relies only on the linear encoding and decoding functionality defined above.
% }}

% \paragrap{Random-band OKVS Construction}
% State-of-the-art OKVSs~\cite{RR22,BPSY23} are based on solving sparse linear system, and this paper exclusively considers the random-band (RB) matrix based construction~\cite{BPSY23}. Let ${\sf H}_1: \mathcal{K}\rightarrow [m]$ and ${\sf H}_2: \mathcal{K} \rightarrow \{0,1\}^w$ for appropriate $w \ll m$ be two random functions. Each key $k\in \mathcal{K}$ is mapped into a sparse vector in $\F^m$ that has only $w$ nonzero components $\mathsf{H}_2(k)$ starting at $\mathsf{H}_1(k)$; this map is written as ${\sf row}: \mathcal{K} \rightarrow \F^m$. The {\sf Encode} and {\sf Decode} algorithms with ${\sf row}$ map are as follows:
% \begin{itemize}
% \item {\sf Encode}($\{(k_i,v_i)\}_{i\in [n]}$): Solve the linear system $$\mat M\cdot \vec p = (v_1, \cdots, v_n)^{\sf T}$$ where $\mat M = [{\sf row}(k_1), \cdots, {\sf row}(k_n)]^{\sf T} \in \F^{n\times m}$, and output the solution $\vec p \in \F^m$. The random-band structure of $\mat M$ enables to solve this linear system in $O(nw)$ time.
% \item {\sf Decode}($k, \vec p$): Output $\langle {\sf row}(k), \vec p \rangle$ over $\F$. The random-band structure reduces this inner product to 
% $$ \sum_{i=1}^w u_i \cdot p_{{\sf H}_1(k) + i - 1} $$ where $\vec u = {\sf H}_2(k)$, which takes $O(w)$ time.
% \end{itemize}
% For an input size $n$ and expansion ratio $\varepsilon$ of $m\approx(1+\varepsilon)n$, the band width $w$ is chosen so that the induced RB matrix $\mat M\in\F^{n\times m}$ has full rank with probability at least \(1-2^{-\lambda}\), where $\lambda = 40$ is a statistical security parameter. This implies the encoding failure probability less than $2^{-\lambda}$.
% The details of the $(\varepsilon,w)$ parameter-selection methodology of~\cite{BPSY23} can be found in \Cref{sec:OKVS experiment}.

%%%%%%%%%%%%%%%기존 OKVS 정의%%%%%%%%%%%%%%%%%
\subsection{Oblivious Key-Value Stores}\label{sec:RB-OKVS}
An oblivious key-value stores~\cite{GPR+21} is a data structure that compactly represents a desired mapping from a set of keys to corresponding values. The definition is as follows:

\begin{definition}
An Oblivious Key-Value Stores (OKVS) is parameterized by a key universe $\mathcal{K}$, a field $\F$ and an expansion ratio $\varepsilon$, and consists of the two algorithms:
\begin{itemize}
    \item $\mathsf{Encode}(T)$:  
    The encode algorithm receives a set of $n$ distinct key-value table 
    $T = \{ (k_1, v_1), \dots, (k_n, v_n) \} \in (\mathcal{K} \times \F)^n$  
    and outputs the encoding $\vec p \in \F^m$ where $m \approx (1 + \varepsilon)\cdot n$.    
    \item $\mathsf{Decode}(k, \vec p)$:  
    The decode algorithm receives the encoding $\vec p \in \F^m$,  
    a key $k \in \mathcal{K}$ and outputs some value $v \in \F$.    
\end{itemize}
\end{definition}

An OKVS is correct if, for any key-value $\{(k_i, v_i)\}_{i\in [n]}$ and $\vec p \leftarrow \mathsf{Encode}(T)$, it holds that ${\sf Decode}(k_i, \vec p) = v_i$ for all $i\in [n]$ with overwhelming probability.

% 우리는 OKVS의 security에 관련된 성질--for example, one might require OKVS encoding to be uniformly random, whenever the value vector is sampled uniformly--을 전혀 쓰지 않기 때문에 이들은 소개하지 않는다.
We do not introduce a separate formal definition of the obliviousness properties of an OKVS, such as the requirement
that the encoding be uniformly distributed whenever the encoded value vector is uniform, because our protocol does
not rely on these properties.
% We nevertheless use the term OKVS throughout the paper because our concrete instantiations are the RB-OKVS construction and its RSB-OKVS transformation.


\paragrap{Random-band OKVS Construction}
State-of-the-art OKVSs~\cite{RR22,BPSY23} are based on solving sparse linear system, and this paper exclusively considers the random-band (RB) matrix based construction~\cite{BPSY23}. Let ${\sf H}_1: \mathcal{K}\rightarrow [m]$ and ${\sf H}_2: \mathcal{K} \rightarrow \{0,1\}^w$ for appropriate $w \ll m$ be two random functions. Each key $k\in \mathcal{K}$ is mapped into a sparse vector in $\F^m$ that has only $w$ nonzero components $\mathsf{H}_2(k)$ starting at $\mathsf{H}_1(k)$; this map is written as ${\sf row}: \mathcal{K} \rightarrow \F^m$. The {\sf Encode} and {\sf Decode} algorithms with ${\sf row}$ map are as follows:
\begin{itemize}
\item {\sf Encode}($\{(k_i,v_i)\}_{i\in [n]}$): Solve the linear system $\mat M\cdot \vec p = (v_1, \cdots, v_n)^{\sf T}$ where $\mat M = [{\sf row}(k_1), \cdots, {\sf row}(k_n)]^{\sf T} \in \F^{n\times m}$, and output the solution $\vec p \in \F^m$. The random-band structure of $\mat M$ enables to solve this linear system in $O(nw)$ time.
\item {\sf Decode}($k, \vec p$): Output $\langle {\sf row}(k), \vec p \rangle$ over $\F$. The random-band structure reduces this inner product to 
$ \sum_{i=1}^w u_i \cdot p_{{\sf H}_1(k) + i - 1} $ where $\vec u = {\sf H}_2(k)$, which takes $O(w)$ time.
\end{itemize}
For an input size $n$ and expansion ratio $\varepsilon$ of $m\approx(1+\varepsilon)n$, the band width $w$ is chosen so that the induced RB matrix $\mat M\in\F^{n\times m}$ has full rank with probability at least \(1-2^{-\lambda}\), where $\lambda = 40$ is a statistical security parameter. This implies the encoding failure probability less than $2^{-\lambda}$.
The details of the $(\varepsilon,w)$ parameter-selection methodology of~\cite{BPSY23} can be found in \Cref{sec:OKVS experiment}.

\begin{remark}
    % 앞서 말했듯 우리에게 필요한 건 Oblivious-KVS가 아닌 correctness만 만족하는 plain key-value store (KVS)이기 때문에, 이를 더 efficient하게 방법으로 구성할 수 있다면 그런 것을 고려해볼만 하다. 하지만 Random-band를 비롯한 최신 OKVS들은 대부분 {\sf Encode}는 linear system solving이며  {\sf Decode}는 inner product 연산이기 때문에 oblivious property들이 필요 없다해도 크게 saving되는 요소가 없다. OKVS라는 용어가 literature에서 많이 쓰이고 있는 바, 굳이 KVS로 구별해서 적지는 않기로 한다.
    Strictly speaking, our protocol requires only the correctness and linear decoding properties of a plain key-value store (KVS), so any more efficient KVS satisfying these requirements could be used.
    However, in modern OKVS constructions including random-band OKVS, encoding is already based on sparse linear-system solving and decoding on a sparse inner product; thus, dropping the obliviousness requirement alone does not offer substantial efficiency savings. We therefore retain the standard term OKVS throughout the paper.
\end{remark}


%%%%%%%%%%%%%%%%
\subsection{\blue{Secret-Shared Private Membership Tests}}
A secret-shared private membership test (ssPMT)~\cite{JSZG24} is a two-party functionality between a sender holding a sequence of distinct elements $Y=(y_1,\cdots,y_{n_y})$ and a receiver holding a set $X$. Its goal is to provide the parties with XOR shares of the membership indicators $\mathbf 1[y_j\in X]$, without revealing the individual
indicators to either party. The formal functionality used in this paper is given in \Cref{fig:F-RL-sspmt}.

\begin{remark}
    % 거의 똑같은 functoinality가 circuit-PSI라고도 자주 불리는데, Circuit-PSI는 보통 associated payload가 존재하는 context에서 사용되는 듯함. 우리는 ssPMT가 좀 더 직관적인 이름이라 생각해서 이걸 선택함.
    Similar functionalities are often referred to as circuit PSI in the literature, especially when the resulting membership shares are consumed by a subsequent secure computation. We use the term ssPMT because it directly emphasizes the output of our primitive: secret shares of per-item membership indicators.
\end{remark}

\paragrap{Meaning of {\sf Reorder}}
% The most simple formulation of ssPMT은 one party의 ordering대로 출력하는 것일 것이다. (\Cref{fig:F-sspmt-informal}이 정확히 그러한 예시임).
% 그러나 정의가 {\sf Reorder}라는 다소 복잡한 개념을 도입한 이유는, state-of-the-art ssPMT 프로토콜들의 output bit share들은 $Y$의 순서를 따르지 않는 경우가 많기 때문이다. 대표적으로 Cuckoo-hashing framework를 따르는 프로토콜들은~\cite{PSTY19,RS21} $Y$를 입력으로 한 cuckoo hash table의 ordering으로 output membership bit를 정의한다: 이 경우 ${\sf Reorder}(Y)$는 cuckoo hashing routine으로 생성되는injective mapping이 될 것.
The simplest formulation of ssPMT outputs the membership shares in the original order of one party's input, as in
\Cref{fig:F-sspmt-informal}. However, the output shares of many efficient ssPMT and circuit-PSI protocols do not follow the original order of $Y$. For example, protocols based on the cuckoo-hashing framework~\cite{PSTY19,RS21} arrange their membership outputs according to the cuckoo table generated from $Y$. In such a protocol, $\mathsf{Reorder}(Y)$ denotes the injective mapping from sender items to their cuckoo-table coordinates.

% Main body에서 제시할 우리의 프로토콜도 이러한 특징을 따르기 때문에, 정의가 다소 지저분하긴 하지만 정확하게 적어두었다.

% 물론 reordering post-process~\cite{CGP20,PRRS24}을 붙여 original $y_i$의 순서대로 다시 output을 재정렬하는 것이 가능하긴 하다. 그러나 그런 방법은 명백히 추가 cost가 필요하다는 단점이 있고, 무엇보다도 our interest인 private set operation을 위해서라면 굳이 그럴 필요가 없다. ${\sf Reorder}(Y)$가 즉 자신의 item $y_i$와 output bit share의 어떤 위치가 대응되는지를 알려주기 때문에 several private set operation들은 충분히 달성할 수 있다, whose expalantion follows below.
% Our protocol has the same general feature: sequencing places
% the sender's query rows in a protocol-defined layout rather
% than in the original order of $Y$. We therefore expose this
% layout explicitly in the ideal functionality. The sender
% cannot choose an arbitrary placement; the map is determined
% by the fixed $\mathsf{Reorder}$ procedure, the sender's input,
% and the public randomness used by the protocol.
Our protocol similarly places sender queries in a protocol-defined layout rather than the original order of $Y$.
We therefore expose this layout through $\mathsf{Reorder}$, which deterministically derives the placement map from the sender's input and public randomness.

% An additional oblivious reordering procedure \cite{CGP20,PRRS24}  could restore the original order of $Y$, but at additional cost. This is unnecessary for our private set operations, since the sender learns $\phi=\mathsf{Reorder}(Y)$ and can directly associate each item (and its payload) with the corresponding output share.

%%%%%%%%%%%%%%%%%%%%%%%%%%%%%%%%%%%%%%%%%%%%%%%%
\begin{figure}[htbp]
\centering
\fbox{%
\begin{minipage}{0.95\linewidth}
\textbf{Parameters:}
Sender $\mathcal S$, receiver $\mathcal R$, set sizes $n_y$ and $n_x$, output length $\ell_{b}$, and a reordering functionality parameterized by $\mathsf{Reorder} \,:\, \mathbb{F}^{n_y} \rightarrow (\phi:[n_y]\rightarrow[\ell_b])\cup\{\bot\}$.\\

\textbf{Input:}
\(\mathcal R\) holds a set \(X\subseteq\{0,1\}^*\). $\mathcal S$ holds a sequence of distinct elements $Y=(y_1,\cdots,y_{n_y})$.\\

\textbf{Functionality:}
\begin{enumerate}
    \item Compute
    $\phi\leftarrow\mathsf{Reorder}(Y)$.
    If $\phi=\bot$, return $\bot$ to both parties.

    \item Initialize $\vec b\gets\vec 0^{\,\ell_b}$.
    
    \item For every $j\in[n_y]$, set $ b_{\phi(j)} \gets \mathbf 1[y_j\in X].$

    \item Sample
    $\vec s\xleftarrow{\$}\{0,1\}^{\ell_b}$.

    \item Return $(\vec s,\phi)$ to $\mathcal S$ and
    $\vec b\oplus\vec s$ to $\mathcal R$.
\end{enumerate}
\end{minipage}
}
\caption{Ideal functionality $\mathcal{F}_{\mathsf{RL\mbox{-}ssPMT}}$.}
\label{fig:F-RL-sspmt}
\end{figure}
%%%%%%%%%%%%%%%%%%%%%%%%%%%%%%%%%%%%%%%%%%%%%%%%%%

\paragrap{Private set operations derived from ssPMT}
It is known that ssPMT (or reverse-PMT) derives several private set operations~\cite{PSTY19,CGS22,CZZ+24}. We summarize them here for completeness. Note that, except PSU, the overhead is about $O(n)$ numbers of oblivious transfer (OT), whose cost is tiny compared to the ssPMT cost.

\son{분량 없으면 아래 appendix로}

First, {\em PSI-Cardinality} that reveals only \(|X\cap Y|\) can be obtained by private evaluation of $\mathsf{HD}(\vec b_\sender, \vec b_\receiver)$, Hamming-distance circuit on the outputs of ssPMT.
As a natural extension, {\em PSI-Threshold} that reveals only \(\vec 1(|X\cap Y| > t)\) comes from additional comparison circuit evaluation, i.e., $\vec 1(\mathsf{HD}(\vec b_\sender, \vec b_\receiver) > t)$.

Assuming the sender has associated payload $d_j$ for each $y_j \in Y$, {\em PSI-Sum} that computes $\sum_{y_j \in X\cap Y} d_j$ can be done as follows.
First, the sender who knows $\phi$ defines \(d_{\phi(j)}'=d_j\) and \(d_k'=0\) for $k \notin \phi(Y)$,
and then two parties engage OT where $\sender$ sets
$(m_0,m_1)=
(r_k + b_{\sender,i}\cdot d_k', r_k + \bar{b}_{\sender, i}\cdot d_k')$,
for a uniform random $r_k$, and $\receiver$ uses \(b_{\receiver, k}\) as its choice bit to have $m_k$.
By each party locally summing \(R=\sum_{k} r_k\) and \(M=\sum_{k}m_k \), it can be checked that $M-R=\sum_{y_j\in X\cap Y}d_j$, which implies PSI-Sum.

{\em Private Set Union} computes $X\cup Y$, or equivalently $Y\setminus X$, as follows.
As a naive approach, following the idea of PSI-Sum, one may set 
\[
(m_0,m_1)=
\begin{cases}
(y_i,\bot) & \text{if } b_{\sender,i}=0,\\[2pt]
(\bot,y_i) & \text{if } b_{\sender,i}=1.
\end{cases}
\]
This satisfies the basic functionality of transmitting only the set difference.
However, it has been repeatedly shown that this approach is insecure due to the placement map $\phi$ determining the ssPMT layout~\cite{KRTW19,JSZG24,LLT24}.
Accordingly, many related works~\cite{GMR+21,CSSW24,JSZG24,TBZ+25,WCS+26} address this issue by obliviously shuffling the layout, formalized as Permute+Share (P\&S)~\cite{CGP20,PRRS24}.
Similarly, PSU can be obtained from ssPMT by first shuffling the membership-bit shares and then applying the above OT-based approach.
For further details, see~\cite{CSSW25}.

%%%%%%%%%%%%%%%%%%%%%%%%%%%%%%%%%%%%%%%%%%%%%%%%%%%%%%%%%%%%%%%
% Sec 3
%%%%%%%%%%%%%%%%%%%%%%%%%%%%%%%%%%%%%%%%%%%%%%%%%%%%%%%%%%%%%%%

\section{RSB Matrices with Homomorphic Decoding}\label{sec:rlwe-okvs}
This section proposes a variant of the Random Band Oblivious Key-Value Stores (RB-OKVS) construction~\cite{BPSY23} that is compatible with RLWE batching, together with a tailored algorithm for efficient homomorphic OKVS decoding. 
Specifically, given encryptions of the OKVS encoding vector $\vec p$, the algorithm homomorphically evaluates $\mat M\vec p$ where $\mat M$ is a plain random-band matrix induced by $K=\{k_i\}_{i\in[n]}$.
More precisely, to obtain the variant RB matrix compatible with batching, \Cref{subsec:RSB OKVS} introduces a matrix family based on \emph{spacing} random bands by the batching slot size $N$. 
Then, \Cref{subsec:sequencing} addresses collisions among starting positions via a sequencing technique that produces multiple matrices, each of which admits an efficient diagonal-block evaluation.
Throughout this section, we assume that the OKVS encoding length $m$ is a multiple of the RLWE slot size $N$, and we write $b = m/N$.

%%%%%%%%%%% Subsec 3.1%%%%%%%%%%%%%%%%%%%%%%%%%%%%
\subsection{Random Spaced Band}\label{subsec:RSB OKVS}
We propose a variant of RB matrix family that places each random band element {\em spaced $N$ positions apart} rather than consecutively, what we call random spaced band (RSB) matrix. From this we expect that, given this variant matrix $\mat M$ is diagonalizable with respect to the starting points, each $N\times m$ sub-matrix of $M$ contains only diagonal entries, so that we only need the ciphertexts $c_i$ that encrypts $\vec p_i$ to compute $\mat M\cdot \vec p$. We now describe this more thoroughly.

\paragrap{Definition of a RSB matrix family}
The $N$-spacing idea is realized by the following column remapping:
\begin{equation}\label{eq:spacing_idea}
1 \mapsto 1, \qquad 2 \mapsto N+1, \qquad 3 \mapsto 2N+1, \qquad \cdots
\end{equation}
However, a naive extension of this rule, namely $i \mapsto (i-1)\cdot N + 1 \pmod m$
is not one-to-one; any index of the form $k\cdot (m/N)+1$ maps to $1$.

To extend the idea in \Cref{eq:spacing_idea} to a column-wise bijection, we consider each column index $i \in [m]$ in its unique representation $i = (i_s - 1)\cdot b + i_b$ where $i_s \in [N],~ i_b \in [b]$ and $b := m/N$, and define a column-index permutation
$\pi_N : [m] \rightarrow [m]$ by $$ 
\pi_N: i = (i_s - 1)\cdot b + i_b   \mapsto   (i_b - 1)\cdot N + i_s.$$
It is immediate to verify that $\pi_N$ is indeed a permutation.

Using this observation, we define the RSB matrix family as the image of the RB matrix family under this fixed column permutation: given any RB matrix $\mat A \in \F^{n\times m}$, its RSB matrix is $\mat B := \mat A \cdot \mat C_N$, where $\mat C_N \in \{0,1\}^{m \times m}$ is the column-permutation matrix induced by $\pi_N$, i.e.,
\begin{equation}\label{eq:N_perm_mat}
  \mat C_N[i][j] =
  \begin{cases}
    1 & \text{if } j = \pi_N(i),\\
    0 & \text{otherwise.}
  \end{cases}
\end{equation}

As observed in \Cref{alg:RSB matrix generato,alg:RSBEncoding}, the RSB matrix generation and OKVS encoding with respect to a RSB matrix can be efficiently done using existing methods for RB matrices~\cite{BPSY23}.


\paragrap{Shape of diagonalized RSB matrices}
We will perform the matrix--vector product after diagonalizing the starting positions of the RSB matrix, and hence we need to understand the resulting placement of elements precisely. Here, however, a minor issue arises: the column permutation $\pi_N$ does not move a consecutive band entirely into an $N$-spaced band. Each time a consecutive band crosses one of the $b$ boundaries, the corresponding spaced band crosses an entire row boundary, and at that point the spaced-band element is placed one position to the right of its natural position modulo $m$. A toy graphical example of this process is provided in \Cref{sec:example}.
% ; see the green points in \Cref{fig:PermutedRSB}, which are shifted by one position from where they would otherwise naturally lie.

As a result, every $N \times m$ submatrix of the diagonalized RSB matrix can have up to $w-1$ non-main diagonals, which in turn implies that we need $w-1$ extra ciphertexts encrypting rotated copies of corresponding parts of $\vec p$; the concrete construction of each such ciphertext is described in detail in \Cref{subsec:okvsdecoding}. Nonetheless, this is still a substantial improvement from the original RB matrix, where a length-$w$ consecutive band requires $w$ slot-rotated copies of $\vec p$, equivalent to $w\cdot b$ ciphertexts.

%%%%%%%%%%% Subsec 3.2%%%%%%%%%%%%%%%%%%%%%%%%%%%%
\subsection{Sequencing} \label{subsec:sequencing} 
Although we defined the RSB matrix family in a form well suited for diagonal matrix--vector product, this approach is effective only when the matrix can be fully diagonalized. However, because $\mathsf{H}_1(k_i)$ is random, collisions among row starting positions are inevitable. %Even a single collision prevents the RSB matrix $\mat{\hat{M}}$ from being perfectly diagonalized, which in turn undermines our diagonal matrix--vector product strategy. 
To address this, we present a technique---which we call \emph{sequencing}---as described below.

\paragrap{Naive sequencing into $m \times m$ matrices} 
A naive remedy is to distribute rows with colliding starting positions across multiple diagonalized $m\times m$ matrices like \Cref{fig:sequencing_overview_b}. 
The number of matrices produced, denoted $L$, equals the maximum number of keys mapped to the same starting position, which can be modeled as the maximum load in the standard balls-into-bins process with $n$ balls and $m$ bins. 
For $m = O(n)$, the maximum load is $O(\log n)$ with high probability (e.g.,~\cite{RS98}), so the naive method produces $O(\log n)$ diagonal matrices and incurs $O(\log n)$ multiplicative computational overhead.

% One might expect the overhead to be smaller because later matrices become sparse as the number of bins occupying more than $i$ items decreases. 
% However, under our diagonalized matrix--vector evaluation, the cost of $\mat{\hat{M}}_i \cdot \vec p$ depends on the number of nonzero $N\times N$ blocks (equivalently, nonzero $N\times m$ row-blocks), since each such block contributes at least one diagonal term in the homomorphic evaluation. 
% Thus, entry-wise sparsity does not translate into a proportional reduction in the number of nonzero blocks, and the overall runtime remains close to $O(\log n)$ times an $m$-dimensional matrix--vector product.

\paragrap{Compact sequencing into $N \times m$ matrices}
We avoid the above overhead by exploiting the structure of the RSB matrix $\mat{\hat{M}}$:
Observe that two rows are aligned to the same diagonal position whenever their starting points are equal modulo $N$.
Thus, we might instead build multiple $N \times m$ matrices by grouping each key $k_i$ into the bin indexed by $\mathsf{H}_1(k_i) \bmod N$. 
This changes the underlying occupancy model from $n$ balls into $m$ bins to $n$ balls into $N$ bins. Since $N\ll n$, the resulting bin loads are considerably more balanced, thereby reducing the wasted sparsity of the full $m\times m$ sequencing approach.

% This enables us to analyze in the $n$-balls-into-$N$-bins model, instead of the $n$-balls-into-$m$-bins model: Since $N \ll n$, the loads become much more balanced, which significantly reduces wasted sparsity of full $m\times m$ sequencing. 
% Moreover, the resulting matrices are well suited for block-level matrix--vector product: the $i$-th row appears only near the $i$-th diagonal within the overall matrix structure.

% The above compact sequencing can be further improved. 
% 하지만, 너무 단순한 $N$ binning은 계산 효율이 떨어진다. For homomorphic evaluation, the computational cost of $\mat M \vec p$ is governed by the number of nonzero $N\times N$ blocks (equivalently, nonzero $N\times m$ row-blocks).
However, naively grouping rows solely according to their starting positions modulo $N$ is still computationally inefficient. The homomorphic cost of evaluating $\mat M\vec p$ is governed by the number of nonzero $N\times N$ blocks, or equivalently by the number of block positions occupied in each $N\times m$ layer.
While the $n$-balls-into-$m$-bins sequencing always yields each $N\times m$ submatrix with exactly $w$ nonzero diagonals, the $n$-balls-into-$N$-bins sequencing scatters row starting points across the entire matrix; it potentially increases the number of nonzero diagonals up to $m/N + w-1$ for every sequenced matrix. 

The natural question is therefore how to reallocate the $n$ rows (with starting points $a_i := {\sf H}_1(k_i)$) into a collection of $N\times m$ matrices, \emph{while keeping the starting points clustered}. 
Regarding, we come up with a greedy sequencing~\Cref{alg:sequencing}. Given a span limit parameter $\zeta$, it inserts each row in ascending index manner as follows: If there is an existing $N\times m$ layer matrix as long as its occupied block span remains at most $\zeta$ after insertion, then put the row in it (line 6 in \Cref{alg:sequencing}); otherwise, create a new $N\times m$ layer. 
It can be easily checked that each resulting $N \times m$ matrix of \Cref{alg:sequencing} has at most $w+\zeta$ cyclic diagonals. One can find a toy graphical example of this process in \Cref{sec:example}.


%%%%%%%%%%%%%%%%%%%%%%%%%%%%
%시퀀싱 알고리즘 수정
\newcommand{\CondLine}[1]{%
  \Statex\hspace{\algorithmicindent}{\quad#1}%
}
\newcommand{\StateBox}[1]{%
  \State \parbox[t]{\dimexpr\linewidth-\algorithmicindent}{#1}%
}

%%%%%%%%%%%%%%%%Sequencing alg%%%%%%%%%%%%%%%%%

\begin{algorithm}[h]
\caption{\textsf{RSB-OKVS.Sequencing}}\label{alg:sequencing}
\begin{algorithmic}[1]
  \Input RSB\mbox{-}matrix $\mat{\hat{M}}$ with $i$-th row starting at $a_i$, and span limit $\zeta$
  \Output Layer matrices $\{\mat{\bar{M}}_j\}_{j=1}^{L}$
\State $L \leftarrow 0$ \InlineComment{Number of layers}

\For{$i \in [n]$, in nondecreasing order of $a_i$}
\State $b_i \leftarrow \lceil\frac{a_i}{N}\rceil$ and $\tau_i \leftarrow a_i~(\bmod~N)$
  \State Find the smallest $\mu \in [L]$ such that 
    % \Indent 
    \CondLine{$- \mat{\bar{M}}_{\mu}[\tau_i]$ is zero 
    \Comment{$\tau_i$-th row is not occupied}}
    \CondLine{$- \max(\beta_\mu,b_i)-\min(\alpha_\mu,b_i) \le \zeta-1$ 
    \Comment{At most $\zeta$ block span} }
     
  \If{such $\mu$ exists} 
    \State Update $(\alpha_\mu,\beta_\mu) \leftarrow (\min(\alpha_\mu,b_i),\max(\beta_\mu,b_i))$
  \Else \InlineComment{Make a new layer matrix}
    \StateBox{Set $L \leftarrow L+1$ and $(\alpha_L,\beta_L)\leftarrow(b_i,b_i)$,\\ 
    and initialize a zero matrix $\mat{\bar{M}}_L \in \Z_p^{N\times m}$}
  \EndIf
  \State $\mat{\bar{M}}_{\mu}[\tau_i] \gets \mat{\hat{M}}[i]$
\EndFor  
  \State \Return $\{\mat{\bar{M}}_j\}_{j=1}^{L}$
\end{algorithmic}
\end{algorithm}
%%%%%%%%%%%%%%%%%%%



Note that the span limit $\zeta$ naturally creates a trade-off between computation and communication cost: A smaller $\zeta$ reduces the number of diagonals $(w + \zeta)$, but increases the resulting number of layers $L$. We choose balancing $\zeta$ for our full protocol, whose choice rationale is presented in \Cref{subsec:RSB parameter choice}.

\paragrap{Optimality of our greedy sequencing}
Surprisingly, \Cref{alg:sequencing} outputs the minimum possible number of layer matrices among all assignments for an input starting positions and block-span limit $\zeta$. Consequently, the number of resulting layers $L$ cannot be reduced merely by replacing our greedy procedure with a better one. Although this optimality result and its proof are of independent algorithmic interest, they are orthogonal to the cryptographic construction. We therefore defer the formal statement and proof to the appendix (\Cref{prop:det-Lprime}).

%%%%%%%%%%%%%%%%%%%%%3.3%%%%%%%%%%%%%%%%%%%%%%
\subsection{Full Homomorphic Decoding Procedure}\label{subsec:okvsdecoding}
The homomorphic decoding algorithm is obtained by combining \Cref{alg:RSB matrix generato} and \Cref{alg:sequencing}, whose full description is given in \Cref{alg:RSBDecoding}. 

We prepare a ciphertext set \(\{c_\alpha\}_{\alpha=1}^{b+w-1}\) where $b=\frac{m}{N}$.
Write $\vec p=(\vec p_1 \,\|\, \cdots \,\|\, \vec p_b)\in\mathbb Z_p^{m(=Nb)}$, where each \(\vec p_i\in\mathbb Z_p^N\) is a contiguous length-\(N\) block.
For $ i\in[b],\ j\ge 0,\ jb+i\le b+w-1$, the ciphertexts form as
\[
  c_{jb+i}=\mathsf{RLWE.Enc}(\rho_j(\vec p_i)),
\]
where \(\rho_j(\vec p_i)\) denote the blockwise cyclic left shift of \(\vec p_i\) by \(j\) slots. 
% (On the first reading, the reader may consider the case where $w \ll b$ case.)
In other words, the first \(b\) ciphertexts encrypt the base blocks \(\vec p_1,\cdots,\vec p_b\), the up to \(w-1\) additional shifted ciphertexts encrypt the corresponding one-slot cyclic shifts, and so on. At the final shift level, fewer than $b$ ciphertexts may remain, corresponding to an initial prefix of the blocks. In particular, a width-\(w\) RSB band starting at block \(i\) uses exactly the consecutive ciphertexts $\{c_i,c_{i+1},\dots,c_{i+w-1}\}$, which are always contained in \(\{c_\alpha\}_{\alpha=1}^{b+w-1}\).

Since each layer matrix $\bar M_i$ obtained by sequencing consists of at most $w + \zeta$ nonzero (super-)diagonals, its encrypted multiplication with $\vec p$ can be carried out with $w + \zeta$ homomorphic slot-wise multiplication and addition, given the above ciphertext set \(\{c_\alpha\}_{\alpha=1}^{b+w-1}\).

%%%%%%%%%%%%%%%%%%%%%%%%%%%%%%%%%%
%Decoding 알고리즘
\begin{algorithm}[h]
\caption{\textsf{RSB-OKVS.HomDecode}}
\label{alg:RSBDecoding}
\begin{algorithmic}[1]
\Input A key set \(K\), a spacing length $N$, a span limit $\zeta$, and ciphertexts 
$\{c_\alpha\}_{\alpha=1}^{b+w-1}$ of the form $c_{jb+i}=\mathsf{RLWE.Enc}(\rho_j(\vec p_i))$ s.t. $jb+i\le b+w-1$, where $\vec p=(\vec p_1 \,\|\, \cdots \,\|\, \vec p_b)$
\Output OKVS decoded ciphertexts \(\{{c}^*_t\}_{t=1}^{L}\)

\State $\mat{\hat{M}} \leftarrow$ \textsf{RSB-OKVS.GenMat}$(K,N)$
\State $\{\mat{\bar M}_i\}_{i\in [L]}\leftarrow$\textsf{RSB-OKVS.Sequencing}$(\mat{\hat{M}}, \zeta)$ 
\For{\(t=1,\cdots,L\)} 
  \StateBox{Compute $c_t^*=\mathsf{RLWE.Enc}(\mat{\bar M}_t\cdot \vec p)$ using $w + \zeta$ {\sf PlainMult} and {\sf CtxtAdd}.}
\EndFor
\State \Return $\{ c_t^*\}_{t=1}^{L}$
\end{algorithmic}
\end{algorithm}
%%%%%%%%%%%%%%%%%%%%%%%%%%%%%%%%%%%%

\paragrap{Reason for pre-rotation, rather than HE-rotation}
Some HE familiar readers may consider letting the sender generate the additional $w-1$ ciphertexts via homomorphic rotation, rather than sending pre-rotated $w-1$ ciphertexts; this will save the communication, with some additional computational cost for homomorphic rotations.
However, our deployed BGV HE~\cite{BGV14} scheme do not directly support a full cyclic rotation over all \(N\) slots. Rather, the packed slots are viewed as a \(2 \times (N/2)\) array, in the sense that native slot manipulations are row/column-wise rotations. This makes the vectorized rotation a bit inefficient, and currently we are not aware of the way to efficiently utilize such structure.


\paragrap{Optimization with full-field random bands}
A key factor of the high performance of RB-OKVS~\cite{BPSY23} was that its band elements are binary $\{0,1\}$. Then, given that the OKVS value field is characteristic-two (e.g., $\F=\mathrm{GF}(2^{\ell})$), the OKVS encoding can be implemented only with row-wise XORs that can be very efficiently performed by modern processors. However, in our setting, we need to work over the prime field $\mathbb{Z}_p$, the plaintext space of RLWE encryption. Then, even though we set the OKVS matrix $\mat M$ as binary, the OKVS encoding must be performed over $\mathbb{Z}_p$, so the base computation unit becomes $\Z_p$ arithmetic, rather than XORs. Thus, binary band elements offer little benefit over $\mathbb{Z}_p$. 

Upon this observation, we sample band elements from the full field $\mathbb{Z}_p$, and this even provides the following benefit. In \cite{BPSY23}, the band-width $w$ (and expansion rate $\varepsilon$) is chosen so the induced RB matrix is full rank with probability at least $1-2^{-\lambda}$.
In this respect, we can expect that the full field bands would increase the entropy of matrix relative to binary bands, and hence make full rank more likely.
In other words, we can expect that the target success probability $1 - 2^{-\lambda}$ can be achieved with a smaller band width $w$ for the full field band.

We indeed observe that, for various $n$ and expansion ratio $\varepsilon$, the same target probability can be attained with smaller band width; the detailed experimental results are presented in \Cref{sec:OKVS experiment}. Since both computational and communication cost of our homomorphic decoding procedure (\Cref{alg:RSBDecoding}) depends on the band width $w$, this directly improves the performance.

%%%%%%%%%%%%%%%%%%%%%  sec5  %%%%%%%%%%%%%%%
\section{\blue{Our Secret-Shared PMT Protocol}} \label{sec:RL-sspmt}
The previous \Cref{sec:rlwe-okvs} established how to efficiently evaluate the RSB-OKVS decoding under RLWE-PKE. 
Building on these encrypted decoding results, one can obtain secret-shared membership indicators by masking each decoded plaintext before decryption and then applying a secret-shared equality test, by following \Cref{fig:zcl_pke_framework}.
% However, the sequencing procedure produces an input-dependent number $L$ of decoding ciphertexts, and returning exactly these $L$ ciphertexts would reveal a function of the sender's set $Y$ through the response length.
Moreover, several security details must be addressed to turn this masking-based idea into a complete protocol. 
The section resolves these issues and presents the resulting protocol, secure against semi-honest adversaries, in \Cref{fig:our-ss-pmt-protocol}.


\subsection{Dummy Layer Padding}\label{subsec:dummy-padding}
The sequencing algorithm produces $L$ layer matrices, and homomorphic batched decoding returns one RLWE ciphertext for each layer. Since the collision pattern of the query-row starting positions is determined by the sender's input $Y$ and the public row-map randomness, the resulting value $L$ may depend on $Y$. Consequently, returning only these $L$ ciphertexts would reveal a function of the sender's input through the response length, regardless of how their contents are masked or rerandomized.

To make the response independent of the sender's input, we fix a public layer budget $L'=L'(n_y,N,m,\zeta)$ such that
\[
    \Pr[L>L']\leq2^{-\lambda},
\]
where the probability is taken over the randomness defining the query-row starting positions. If $L>L'$, the protocol returns $\bot$; otherwise, the sender pads the response and returns exactly $L'$ ciphertexts. The following lemma gives a concrete method for selecting this layer budget.

\providecommand{\Bin}{\mathrm{Bin}}
\providecommand{\ext}{\mathrm{ext}}

\begin{lemma}[Layer budget] \label{lemma:Lprime}
    Let $N \mid m$, $b = m/N$, $\zeta \ge 1$, and let the starting points $a_1, \dots, a_n$ be drawn independently and uniformly from $[m]$. For $g \in [b]$ define the binomial cap $$ C(g) \;=\; \min\Big\{ k \in \mathbb{N} \;:\;  \Pr\big[\, \Bin(n,\, g/m) \ge k \,\big] \;\le\; \frac{2^{-\lambda}}{N b^2} \Big\},$$ and set $ L'(n, N, m, \zeta) \;=\; \max_{1 \le g \le b} \left\lceil \frac{\big(C(g)+1\big)\,(b+\zeta-1)}{g+\zeta-1} \right\rceil .$ \\
    Then the number $L$ of matrices output by \Cref{alg:sequencing} on $\{a_i\}_{i \in [n]}$ with span limit $\zeta$ satisfies $\Pr\big[\, L > L' \,\big] \;\le\; 2^{-\lambda}.$ 
    \\ Furthermore, if $\zeta \ge \min(\lambda, b)$ and $\lambda \ge \log_2(Nb^2)$,
then $L' = O(n/N+\lambda)$.
\end{lemma}

\begin{proof}
    Can be found at \Cref{app:Lprime}.
\end{proof}

\begin{remark} 
    An similar issue in the cuckoo/simple-hashing framework for unbalanced PSI~\cite{CLR17} that also resolves it with padding, and later works~\cite{CHLR18,CMdG+21} eliminate such concern using OPRF preprocessing. However, the OPRF solution only applies for their PSI case, not our ss-PMT setting, because OPRF can only randomize non-intersecting elements. %In PSI where the intersectin element  but violates our requirement that no per-item membership information be disclosed. 
\end{remark}

\paragrap{Choice of padding ciphertexts}
For correctness, the padded ciphertexts must ensure that every slot in the additional $L'-L$ layers has an underlying value of zero before masking.
For every $\mu\in\{L+1,\cdots,L'\}$, the sender therefore generates a fresh encryption of the all-zero vector and adds an independent mask $\vec r_\mu\xleftarrow{\$}\mathbb Z_p^N$.
The receiver consequently decrypts the padded ciphertext to $\vec d_\mu=\vec r_\mu$. Since $\mathsf{IS}\in\mathbb Z_p^\ast$, the subsequent equality test compares $r_{\mu,i}$ with $r_{\mu,i}-\mathsf{IS}$ and thus deterministically outputs zero at every padded coordinate.

% \paragrap{Asymptotics on $L'$ and concrete choice of $\zeta$}
% {\red{We can now settle the trade-off left open in \Cref{subsec:sequencing}. Increasing $\zeta$ decreases the number of layers---both the realized $L$ and the budget $L'$ of \Cref{lemma:Lprime}---but inflates the per-layer decoding work of $w+\zeta$ diagonals; recall that the return communication and the number of equality tests are driven by $L'$, while the homomorphic decoding cost is driven by $L$. The decisive observation is how $L'(\zeta)$ decays: it falls steeply at first, then quickly saturates at an information-theoretic floor, namely the maximum bin load of the $n$ starting points (rows of one bin always need pairwise-distinct layers), whose $2^{-\lambda}$ quantile no budget can undercut. At $n=2^{20}$, for instance, $L'(\zeta)$ decreases as $530,\, 359,\, 277,\, 252$ for $\zeta = 10,\, 20,\, 40,\, 74$, against a floor of $250$. We therefore grow $\zeta$ exactly until the decrease of $L'$ stalls, and stop there; the resulting values are reported in \Cref{tab:rsb_params}, where $L'$ sits within $2$--$10\%$ of the floor.\footnote{The implementation automates this rule: it scans $\zeta$, weighing the fixed cost of a layer (one returned ciphertext plus $N$ equality tests) against that of a diagonal (one plaintext multiplication) with a single measured machine constant, and among all $\zeta$ within $1\%$ of the optimum it returns the largest---communication only improves with $\zeta$, so the far edge is preferable.}}}

\paragrap{Overheads on performance}
Fixed-length padding increases the number of returned ciphertexts to $L'$, and hence increases the return communication, the receiver's decryption cost, and the number of secret-shared equality tests to $NL'$. In contrast, the main computational bottleneck, namely homomorphic OKVS decoding, is evaluated only for the actual $L$ layers and is therefore unaffected by the additional $L'-L$ padding layers.

\paragrap{Concrete behavior of $L'$ and choice of $\zeta$}
Increasing $\zeta$ decreases both the sequenced results $L$ and the budget $L'$, but inflates the per-layer decoding work of $w+\zeta$ diagonals. 
As a decisive observation, $L'$ has a definite lower bound--the maximum bin load--, and it rapidly falls to that bound as $\zeta$ grows; 
for example, at $n=2^{20}$, $L'$ decreases as $530,\, 359,\, 277,\, 252$ towards its lower bound $250$ for $\zeta = 10,\, 20,\, 40,\, 74$. 
Our concrete choice of $\zeta$ is some proper point where the decrease of $L'$ stalls.
The resulting values are reported in \Cref{tab:rsb_params}.


%%%%%%%%%%%%%5.1%%%%%%%%%%%%%%
\subsection{Full Protocol}\label{subsec:ourRLsspmtProtocol}

%%%%%%%%%%%%%%%%%%%%%%%%%%%%%%%%%%%%%%%%%%%%%%%%%%
%또 수정한 프로토콜
\begin{figure*}[htbp]
\centering
\fbox{%
\begin{minipage}{0.95\textwidth}

\textbf{Input.}
The sender \(\sender\) holds \(Y=\{y_j\}_{j\in[n_y]}\), and the receiver \(\receiver\) holds \(X=\{x_i\}_{i\in[n_x]}\). The output length $\ell_b=NL'$.

\textbf{Output.}
$\sender$ obtains  $\vec b_Y=(b_{Y,k})_{k\in[\ell_b]}$ together with a map $\phi:[n_y]\rightarrow[\ell_b]$, while  $\receiver$ obtains  $\vec b_X=(b_{X,k})_{k\in[\ell_b]}$.

\medskip
\noindent
\textbf{1) Setup.}

\begin{itemize}
    \item The parties agree on RLWE parameters $(N,q)$, plaintext modulus $\log p \ge \lceil \lambda + \log n_y \rceil$, a sequencing span parameter $\zeta$, a fixed layer budget $L'$, and then RSB\mbox{-}OKVS parameters $w$, $m$ s.t. $b:= m/N \in \Z$.
    \item $\receiver$ generates a RLWE-PKE key pair $(\mathsf{pk}, \mathsf{sk})$, and sends $\mathsf{pk}$ to $\sender$.
    % \item  The parties generate all input-independent correlated randomness required for \(\ell_b=NL'\) invocations of $\mathcal{F}_{\mathsf{ss\mbox{-}peqt}}$.
\end{itemize}

\medskip
\noindent
\textbf{2) RSB-OKVS encoding and encryption.}

\begin{itemize}
    \item  $\receiver$ samples a nonzero indicator \(\mathsf{IS}\xleftarrow{\$}\mathbb Z_p^\ast\) and computes $\vec p\leftarrow \mathsf{RSB\mbox{-}OKVS.Encode}  \bigl(\{(x_i,\mathsf{IS})\}_{i\in[n_x]}\bigr) \in\mathbb Z_p^m$.
    \item $\receiver$ partitions \(\vec p=(\vec p_1\|\cdots\|\vec p_b)\), where \(\vec p_i\in\mathbb Z_p^N\), and generates \(\mathcal C=\{c_\alpha\}_{\alpha=1}^{b+w-1}\), where $c_{jb+i} \leftarrow \mathsf{RLWE.Enc}_{\mathsf{pk}} \bigl(\rho_j(\vec p_i)\bigr)$  for every \(i\in[b]\) and \(j\geq 0\) satisfying  \(jb+i\leq b+w-1\).
    \item $\receiver$ sends \(\mathcal C=\{c_\alpha\}_{\alpha=1}^{b+w-1}\) to $\sender$.
\end{itemize}

\medskip
\noindent
\textbf{3) Homomorphic RSB-OKVS decoding, masking, and decryption.}

\begin{itemize}
    \item $\sender$ computes $ \{c_\mu^*\}_{\mu=1}^{L}  \leftarrow \mathsf{RSB\mbox{-}OKVS.HomDecode} (Y,N,\zeta,\mathcal C)$ (\Cref{alg:RSBDecoding}), and pads it with dummy $L' - L$ encryptions of zero, precisely $c_\mu^* \leftarrow \mathsf{RLWE.Enc}_{\mathsf{pk}}(\vec 0)$. $\sender$ also records the placement map $\phi:[n_y]\rightarrow[\ell_b]$.
    
    \item For \(\mu\in[L']\), $\sender$ independently samples $ \vec r_\mu \xleftarrow{\$}\mathbb Z_p^N$ and masks each ciphertext by $c_\mu \leftarrow \mathsf{RLWE.PlainAdd}(c_\mu^*,\vec r_\mu)$.

    % \item $\sender$ re-randomizes and masks the \(L\) homomorphically decoded ciphertexts and pads the resulting sequence to the fixed length \(L'\) as follows:
    % \begin{enumerate}
    %      \item For \(\mu\in[L]\), $\sender$ computes $ \bar c_\mu \leftarrow  \mathsf{RLWE.CtxtAdd}\bigl( \mathsf{RLWE.PlainAdd}(c_\mu^*,\vec r_\mu), c_\mu^0 \bigr).$

    %     \item For  \(\mu\in\{L+1,\cdots,L'\}\), $\sender$ generates a masked encryption of the all-zero decoded vector as $\bar c_\mu \leftarrow \mathsf{RLWE.PlainAdd}(c_\mu^0,\vec r_\mu)$.
    % \end{enumerate}
    
   
    \item $\sender$ re-randomizes \(\{c_\mu\}_{\mu\in[L']}\){\footnote{$\sender$ re-randomizes ciphertext, by adding a fresh encryption of zero with flooding noise sufficient to achieve statistical circuit privacy.}} and send to $\receiver$, who decrypts them to have $\vec d_\mu=(d_{\mu,i})_{i\in[N]}$.%  = \mathsf{RLWE.Dec}_{\mathsf{sk}}(\bar c_\mu) \in\mathbb Z_p^N.$ 
\end{itemize}

\medskip
\noindent
\textbf{4) Secret-shared equality testing.}

\begin{itemize}
    \item Two parties flatten $\{\vec r_\mu\}_{\mu \in L'}$ and $\{\vec d_\mu\}_{\mu \in L'}$ to $NL'$-length vector, say $\{r_k\}_{k\in [NL']}$ and $\{d_k\}_{k \in [NL']}$. 
    
    \item For every $k \in [NL']$, two parties invoke $\mathcal{F}_{\mathsf{ss\mbox{-}peqt}}$ on inputs \(r_{k}\) and \(d_{k}-\mathsf{IS}\), and obtain \(b_{Y,k}\) and \(b_{X,k}\), respectively.

    \item $\sender$ returns $\{b_{Y,k}\}_{k\in [NL']}$ and $\receiver$ returns $\{b_{X,k}\}_{k\in [NL']}$.
\end{itemize}

\end{minipage}
}
\caption{Our ssPMT protocol}
\label{fig:our-ss-pmt-protocol}
\end{figure*}
% %%%%%%%%%%%%%%%%%%%%%%%%%%%%%%%%%%%%%%%%%%%%%%%%%%

{\blue{
\paragrap{Sequencing-induced placement map}
The sequencing algorithm (\Cref{alg:sequencing}) first induces an injective placement map $\widetilde{\phi}:[n_y]\rightarrow[NL].$
More precisely, if the query row $\mathsf{row}(y_j)$ is placed in the $i_j$-th row of layer $\mu_j\in[L]$, then $\widetilde{\phi}(j):=(\mu_j-1)N+i_j.$
Since the protocol pads the decoding results to the fixed length $L'$ when $L \leq L'$, it naturally embeds the sequencing output domain $[NL]$ into the fixed output domain $[\ell_b]=[NL'].$
Accordingly, we extend the sequencing-induced map to
$\phi:[n_y]\rightarrow[\ell_b]$ by setting
\[
    \phi(j):=\widetilde{\phi}(j)
            =(\mu_j-1)N+i_j.
\]
Thus, $\operatorname{Im}(\phi)\subseteq[NL]\subseteq[\ell_b]$.
The coordinates outside $\operatorname{Im}(\phi)$ correspond either to unoccupied rows in the first $L$ layers or to slots in the additional $L'-L$ padded layers; their underlying decoding values before masking are all zero. The sender retains $\phi$ locally and outputs it together with its membership-share vector, while the receiver does not learn $\phi$. 

% \paragrap{Sequencing-induced placement map}
% sender는 시퀀싱 알고리즘 수행 및 $L'$개의 디코딩 결과 암호문 전달을 위해 패딩하는 과정을 반영하여 시퀀싱 알고리즘 출력 map을 $[n_y] \rightarrow [\ell_b]$ where $\ell_b =NL'$로 확장시킨다. 
% More precisely, \Cref{alg:sequencing} 과정에서, every query row $\mathsf{row}(y_j)$ is assigned to exactly one row $i_j\in[N]$ of one actual layer $\mu_j\in[L]$. The sender records the corresponding placement map
% \[
%     \phi(j)
%     :=
%     \mathsf{idx}(\mu_j,i_j)
%     =
%     (\mu_j-1)N+i_j.
% \]
% Since each sender query is placed exactly once and each row of a layer contains at most one query, $\phi$ is injective.
% The sender retains $\phi$ locally and does not reveal it to the receiver.

% The coordinates outside $\operatorname{Im}(\phi)$ consist of unoccupied rows in the first $L$ layers and all slots in the
% additional $L'-L$ padded layers. An unoccupied row is the all-zero row of its sequencing matrix, while every padded
% ciphertext encrypts the all-zero vector before masking.

% Hence, letting $v_k$ denote the underlying unmasked decoding value at coordinate $k\in[\ell_b]$, we have
% \[
%     v_k
%     =
%     \begin{cases}
%       \mathsf{Decode}(y_j,\widehat{\vec p}),
%         & \text{if }k=\phi(j)\text{ for some }j\in[n_y],\\
%       0, & \text{if }k\notin\operatorname{Im}(\phi).
%     \end{cases}
% \]
% Consequently, since $\mathsf{IS}\in\mathbb Z_p^\ast$, the secret-shared equality tests produce
% \[
%     b_{Y,k}\oplus b_{X,k}
%     =
%     \begin{cases}
%       \mathbf{1}[y_j\in X],
%         & \text{if }k=\phi(j)\text{ for some }j\in[n_y],\\
%       0, & \text{if }k\notin\operatorname{Im}(\phi).
%     \end{cases}
% \]
}}


\paragrap{Re-randomization and noise flooding}
Although the uniform masks hide the underlying decoding plaintexts, ciphertext rerandomization is still necessary.
Indeed, the receiver holds the RLWE secret key and may otherwise distinguish evaluation-dependent information remaining in the returned ciphertexts, including their noise distributions. The return procedure must therefore provide statistical circuit privacy, rather than merely plaintext privacy.

For every actual and padded layer, the sender adds an independent encryption of zero whose noise includes a sufficient flooding margin. The flooding distribution is chosen so that the resulting ciphertext distribution is statistically independent of the homomorphic decoding path and of whether the ciphertext corresponds to an actual or a padded layer. Unlike ElGamal rerandomization, where adding a fresh encryption of zero suffices, RLWE-based rerandomization must also statistically hide the circuit-dependent noise.

This flooding increases the noise magnitude and therefore requires a sufficiently large return-level ciphertext modulus $q$ to preserve correct decryption. Importantly, it is not sufficient for actual and padded ciphertexts merely to have similar nominal noise magnitudes; their returned distributions must be statistically close under the claimed circuit-privacy guarantee.

%%%%%%%%%%%%%%%%%%%%%%%%%%%%%%%%%%%%%%%%%%%%%%%%%%%%

\paragrap{Parameter setting for correctness}
First of all, note that too short indicator bit-length $\ell_{\sf IS}$ incurs a false positive; a decoded value for an element $y_i \notin X$ may match ${\sf IS}$.
For an element $y_i\notin X$, the full-field band coefficients $H_2(y_i)\in\mathbb Z_p^w$ are uniform and independent of the encoding, since $H_2$ is modeled as a random oracle.
Therefore, the decoded value is uniform over $\mathbb Z_p$. 
Since we set $\log p \ge \ell_{\sf IS}$, applying a union bound over all $y_i \in Y$ motivates setting $\ell_{\sf IS} = \lambda + \lceil \log n_y \rceil$, so that $\Pr[\mathsf{false\_match}] \le 2^{-\lambda}$.
One might doubt the union bound should be taken on every $NL'$ slot. However, as the unoccupied rows and every padded row will be the pure zero row,  its decoding also becomes zero, and thus never match the nonzero indicator $\mathsf{IS}$.

% \magenta{Recall that padded ciphertext 부분은 절대 false positive를 유발하지 않기 때문에, 그 부분은 union bound에 넣지 않아도 된다.}
% For a single element, the event occurs with probability $2^{-\ell_{\sf IS}}$, applying the union bound over all $y_i \in Y$ motivates setting $\ell_{\sf IS} = \lambda + \lceil \log n_y \rceil$, so that $\Pr[\mathsf{false\_match}] \le 2^{-\lambda}$. 
% {\blue{For a single element $y_i \notin X$, the event occurs with probability $2^{-\ell_{\sf IS}}$, since a random decoding value lies in $\mathbb{Z}_p$ and we set $\log p \ge \ell_{\sf IS}$. Applying a union bound over all $y_i \in Y$ motivates setting $\ell_{\sf IS} = \lambda + \lceil \log n_y \rceil$, so that $\Pr[\mathsf{false\_match}] \le 2^{-\lambda}$.}}
The OKVS value field is $\ell_{\sf IS}$ bits, so we require RLWE plaintext modulus satisfying $\log p \ge \ell_{\sf IS}$.
We choose the RLWE ciphertext modulus $q$ to support depth-1 ${\sf PlainMult}$, namely $\log q \gtrsim 2 \cdot \log p$.\footnote{Evaluating a depth-$d$ circuit typically requires $\log q \approx (d+1)\cdot \log p$.}
Finally, we select the batching slot size (equivalently, the RLWE ring dimension) $N$ according to standard RLWE security guidelines~\cite{HEStandard} for the chosen modulus $q$ and the target security level (e.g., 128-bit security). Concretely, we use $N = 2^{13}$ for all set sizes $n$.

\paragrap{Optimizations from RLWE details}
%In a standard RLWE\mbox{-}based encryption scheme, a ciphertext \(c=(c_0, c_1) \in R_q^2\) consists of two polynomials.%, so its bit\mbox{-}length is about $2N \log q$ bits for ring dimension $N$. 
We apply two optimizations from RLWE literature, which reduces the communication almost half. These optimizations are not described in \Cref{fig:our-ss-pmt-protocol} for the sake of simplicity, but indeed applied for experimental results in \Cref{sec:performance}.

First, when the fresh ciphertexts of OKVS encoding, the second component $c_1$ of each RLWE ciphertext $(c_0,c_1) \in R_q^2$ can be compressed into with a short seed, so that the sender can locally expand this seed to reconstruct $c_1$.

Second, for the returning decoded ciphertexts, %Instead, as no further homomorphic computation is needed, 
the receiver can switch down the ciphertext modulus. Combined with noise flooding technique, the concrete returning modulus $\log q'$ is set by $\approx \log p + \log b + \log \lambda$ bit where $\log b$ is the protocol-derived noise and $\log \lambda$ is for noise flooding.
Concretely, we observe that $\log b \approx 10$, respectively. \magenta{The corresponding return-level modulus sizes are $108,112,116,120$ bits, respectively.}

% \blue{Concretely, for $n_y=2^{16},2^{18},2^{20},2^{22}$, we use $\ell_{\sf pre}=8,10,12,12$, respectively, and thus set $\ell_{\sf flood}=48,50,52,52$. The corresponding return-level modulus sizes are $108,112,116,120$ bits, leaving nominal post-flooding decryption margins of $3,3,3,7$ bits, respectively.}


{\blue { 
\begin{theorem} \label{thm:semi-honest-rl-sspmt}
     For $\ell_b=NL'$, the protocol of \Cref{fig:our-ss-pmt-protocol} securely realizes \Cref{fig:F-RL-sspmt} against semi-honest adversaries assuming IND-CPA security of the underlying RLWE encryption.
\end{theorem}

\begin{proof}
    appendix or sketch?
\end{proof}

}}

\paragrap{Asymptotic cost analysis}
Our ssPMT protocol (\Cref{fig:our-ss-pmt-protocol}) has $O(n)$
computation and communication complexity in the input size $n$ (where
$n=n_x=n_y$). Due to the space limit, rigorous details can be found in
\Cref{sec:asymp_cost_anal}, and we briefly explain the idea. The main
reason is that RLWE batching packs $N$ values into one ciphertext: the
receiver sends $b+w-1 = O(n/N+\lambda)$ ciphertexts, the sender returns
$L' = O(n/N+\lambda)$ ciphertexts by \Cref{lemma:Lprime}, and each
ciphertext has size $O(N)$. The $NL'$ equality tests likewise cost
$O(n)$. Moreover, homomorphic decoding performs $L \le L'$ layer
computations, each using $w+\zeta$ slot-wise operations, which is
$O(\lambda)$ per layer since $w = O(\lambda)$ and our choices of $\zeta$
exceed $\min(\lambda, b)$ by at most a small constant factor. Thus, for
fixed security parameters, both communication and computation are linear
in $n$.

%%%%%%%%%%%%%%%%%%%%%%%%%%%%%%%%%%%%%%%%%%%%%%%%%%%%%%%%%%%%%%%%%%%%%%%%%

%%%%%%%%%%%%%%%%%%%%%%%%%%%%Sec 6%%%%%%%%%%%%%%%%%%%%%%%%%%%%%%%%%%%%%%%

\section{Experimental Results}\label{sec:performance}
In this section, we report experimental results for our ssPMT protocol, which serves as the main building block for several private set operations, and for our PSU protocol. We compare their performance with that of state-of-the-art protocols. {\blue{공간이 부족해서 비교하지는 않지만 우리 구현에는 PSI-cardinality, Threshold 등 PSO 프로토콜도 구현되어있다.(언급할까요?)}} All parameters are selected to provide $\lambda=40$ bits of statistical security and $\kappa=128$ bits of computational security, as commonly adopted in the literature. 

\paragrap{Experiment environments}
% {\color{blue}For a fair comparison between our RSB-PSU and existing protocols, we deployed and evaluated their open-source implementations on our local server. }\footnote{\cite{JSZG24}:\texttt{https://github.com/yanxue820/SecurePSU},
% \cite{TBZ+25}:\texttt{https://zenodo.org/records/14725816},
% \cite{ZCL+24}:\texttt{https://github.com/alibaba-edu/mpc4j?tab=readme-ov-file},
% \cite{CZZ+24}:\texttt{https://github.com/yuchen1024/Kunlun/tree/master/mpc}}
Our experiment is conducted on a workstation equipped with AMD Ryzen Threadripper 9970X (4.0GHz) CPU and 64GB RAM. The network environment is simulated by localhost on the same machine while the traffic is controlled by the Linux \texttt{tc} command. Each party runs over a single thread, which is the same setting to the previous works.

\paragrap{Implementation details}
Our implementation in C++ is available at \cite{RLWE-OKVS-git}. For the underlying RLWE-based AHE primitive, we rely on SEAL library~\cite{SEAL} with HEXL~\cite{IntelHEXL} AVX-512 acceleration. 
We also adapt the public RB-OKVS implementation~\cite{RB-github} to implement our RSB-OKVS. Finally, we rely on \texttt{libOTe}~\cite{libOTe} for OT-extension backend.

%%%%%%%%%%%%%%%%%%%%%%%%6.1%%%%%%%%%%%%%%%%%
\subsection{Concrete Parameter Choices} \label{subsec:RSB parameter choice}
Our ssPMT protocol involves several parameters. We explain the concrete choices summarized in \Cref{tab:rsb_params}. 
\kwag{$L'$ 열 추가 및 파라미터 업데이트하기 (준준)} \kwon{완료}

\begin{table}[h]
    \centering
    \caption{Our deployed parameters for RSB-OKVS: $\varepsilon$ is an OKVS
    row expansion ratio, $w$ is a band width, $\zeta$ is a span limit for
    the matrix sequencing. $L$ is the mean of the number of sequenced
    layers, measured by taking average over 100 independent random
    instances. $L'$ is the layer budget of \Cref{lemma:Lprime}, the number
    of ciphertexts the sender actually returns, computed from the public
    parameters alone; see \Cref{subsec:dummy-padding} for the choice of
    $\zeta$.}
    \label{tab:rsb_params}
    \setlength{\tabcolsep}{6pt}

    \begin{tabular}{c|c|ccc|cc|c}
        \toprule
        $N$ & $n$ & $\varepsilon$ & $w$ & $\zeta$ & $L$ & $L'$ & $n/N$ \\
        \midrule
        \multirow{4}{*}{$2^{13}$}
        & $2^{16}$ & 1.1 & 29 & 17 & 20.7  & 47  & 8   \\
        & $2^{18}$ & 0.7 & 43 & 40 & 56.0  & 99  & 32  \\
        & $2^{20}$ & 0.7 & 44 & 52 & 172.8 & 262 & 128 \\
        & $2^{22}$ & 0.7 & 46 & 96 & 602.6 & 789 & 512 \\
        \bottomrule
    \end{tabular}
\end{table}

\magenta{First, recall the batching slot size, namely the RLWE ring dimension \(N\), is fixed to \(N=2^{13}\) in all experiments, which is the smallest choice that securely supports our homomorphic OKVS decoding procedure. We then consider the band expansion ratio \(\varepsilon\) and the band width \(w\). Recall that the requirement for \((\varepsilon,w)\) for our RSB matrix is exactly the same as in the original RB matrix~\cite{BPSY23}: One should choose them so that the encoding failure probability is \(2^{-\lambda}\). However, once the additional features and optimizations specific to our protocol are taken into account, the parameter pairs \((\varepsilon,w)\) suggested in \cite{BPSY23} are no longer optimal in our setting. Therefore, while we follow the methodology of \cite{BPSY23}, we select \((\varepsilon,w)\) differently from their choices.

Basically, the choice of \((\varepsilon,w)\) involves the following trade-off. Since the OKVS encoding length is \(m\approx(1+\varepsilon)n\), increasing \(\varepsilon\) directly increases the communication cost for the fresh ciphertexts sending from $\receiver$ to $\sender$. On the other hand, as the rows become longer, the width \(w\) required for the R(S)B matrix to be full rank becomes smaller, which reduces the computational cost. However, in our protocol, the number of fresh ciphertexts is \(m/N + w - 1\). That is, the communication cost is also affected by \(w\). Consequently, up to a certain point, increasing \(\varepsilon\) can be expected to reduce both communication and computation at the same time. The parameters in \Cref{tab:rsb_params} are chosen as balancing one, from internal experiments. Naturally, different parameter choices may be optimal in different environments, and some use cases may require a different computation--communication trade-off; we provide relevant data in \Cref{sec:OKVS experiment}.

Lastly, the span limit $\zeta$ is chosen per input size as described in
\Cref{subsec:dummy-padding}. Two consequences are visible in
\Cref{tab:rsb_params}. First, $L'$ stays close to its lower bound, the
maximum bin load ($47$, $99$, $262$, $789$ against floors of $46$,
$98$, $250$, $749$). Second, the padding overhead $L'/L$ decreases as
$n$ grows, from $2.3\times$ at $n=2^{16}$ to $1.3\times$ at $n=2^{22}$,
so the relative cost of the fixed-length padding shrinks at scale.
}

%%%%%%%%%%%%%%%%%%%%%%%%6.2%%%%%%%%%%%%%%%%
\subsection{Performance Evaluations}
This section provides a concrete performance evaluation, and comparison with recent PSU and RPMT/ssPMT protocols where all numbers are obtained from running the public implementations%,\footnote{
% \cite{CZZ+24}: \texttt{https://github.com/yuchen1024/Kunlun} \\
% \cite{PT26}: \texttt{https://github.com/asu-crypto/IBLT-based-PSU}\\
% \cite{JSZG24}: \texttt{https://github.com/yanxue820/SecurePSU.git}\\
% \cite{TBZ+25}:\texttt{https://zenodo.org/records/14725816}\\
% \cite{WCS+26}: \texttt{https://github.com/sabakioukenasai/sPSO.git}} 
on our machine. We consider four network environments: 10Gbps with 0.1 ms latency, and 1Gbps, 100Mbps, 10Mbps bandwidth with 80 ms latency. The first setting will be referred as LAN network, and the others are as WAN networks.

%%%%%%%%%%%%%%%%%%%%%%%%PSO%%%%%%%%%%%%%%%%%%%%%%%%%%%%%%%%%%%%%
\paragrap{Comparison of PSO base}
\Cref{tab:sspmt-comparison} compares the RPMT of \cite{CZZ+24} with the ssPMT protocols of \cite{WCS+26} and ours, which serve as core building blocks for PSO. \cite{CZZ+24} performs costly X25519-based cwPRF evaluations for every input item, whereas our RLWE construction amortizes ciphertext-level operations over \(N\) packed slots. Our protocol is therefore faster in every LAN experiment, with the advantage reaching \(3.65\times\) at \(n=2^{22}\). \cite{CZZ+24}, however, requires substantially less communication and is thus generally faster in bandwidth-constrained WAN settings, particularly for smaller sets and at \(10\) Mbps.  More importantly, \cite{CZZ+24} does not prevent DEL, whereas our ssPMT reveals only secret shares of the membership bits.  

\cite{WCS+26} also prevents DEL, but its ssPMT construction relies on communication-intensive secure equality testing and permutation-correlation. These components lead to substantially larger end-to-end communication and computation. Across the reported set sizes, \cite{WCS+26} requires \(6.66\)--\(23.17\times\) the communication of ours, causing its performance to degrade sharply as bandwidth decreases. Consequently, our total runtime is lower in every reported setting, achieving up to a \(23.11\times\) speedup at \(n=2^{22}\) and \(10\) Mbps (\(13077.83\) versus \(565.84\) seconds). For example, at \(n=2^{20}\) and \(100\) Mbps, our protocol achieves a \(14.42\times\) speedup (\(393.04\) versus \(27.26\) seconds).

% =========================================================
% 1. sspmt comparison table(revised_v2)
% =========================================================
\begin{table}[t]
  \centering
  \caption{Comparison of PSO base protocols. DEL-leaky base protocol (RPMT)~\cite{CZZ+24} is reported as baseline, so excluded from colored ranking.}
  \label{tab:sspmt-comparison}

  \footnotesize
  \renewcommand{\arraystretch}{1.15}

  \resizebox{\columnwidth}{!}{%
  \begin{tabular}{ccccccc}
    \toprule
    \multirow{2}{*}{\makecell[c]{Set\\Size}} &
    \multirow{2}{*}{Protocol} &
    \multirow{2}{*}{\makecell{Comm.\\(MB)}} &
    \multicolumn{4}{c}{Runtime (s)} \\
    \cline{4-7}
    & & & 10Gbps & 1Gbps & 100Mbps & 10Mbps \\
    \midrule

    \multirow{3}{*}{$2^{16}$}
      & \cite{CZZ+24}
      & 4.67
      & 2.45
      & 3.31
      & 3.46
      & 6.37 \\

    \arrayrulecolor{gray!35}
    \cmidrule{2-7}
    \arrayrulecolor{black}

      & \cite{WCS+26}
      & 234.44
      & 2.57
      & 13.47
      & 28.64
      & 206.21 \\

      & {\bf Ours}
      & \BEST{35.21}
      & \BEST{1.81}
      & \BEST{4.40}
      & \BEST{6.15}
      & \BEST{33.42} \\

    \midrule

    \multirow{3}{*}{$2^{18}$}
      & \cite{CZZ+24}
      & 18.66
      & 9.80
      & 11.12
      & 11.95
      & 23.99 \\

    \arrayrulecolor{gray!35}
    \cmidrule{2-7}
    \arrayrulecolor{black}

      & \cite{WCS+26}
      & 933.71
      & 10.17
      & 39.76
      & 101.44
      & 817.64 \\

      & {\bf Ours}
      & \BEST{75.45}
      & \BEST{4.41}
      & \BEST{7.88}
      & \BEST{11.60}
      & \BEST{70.24} \\

    \midrule

    \multirow{3}{*}{$2^{20}$}
      & \cite{CZZ+24}
      & 74.66
      & 39.20
      & 42.14
      & 46.59
      & 99.18 \\

    \arrayrulecolor{gray!35}
    \cmidrule{2-7}
    \arrayrulecolor{black}

      & \cite{WCS+26}
      & 3731.01
      & 40.52
      & 146.69
      & 393.04
      & 2019.57 \\

      & {\bf Ours}
      & \BEST{199.17}
      & \BEST{13.57}
      & \BEST{18.96}
      & \BEST{27.26}
      & \BEST{182.65} \\

    \midrule

    \multirow{3}{*}{$2^{22}$}
      & \cite{CZZ+24}
      & 298.63
      & 158.97
      & 168.52
      & 187.18
      & 404.55 \\

    \arrayrulecolor{gray!35}
    \cmidrule{2-7}
    \arrayrulecolor{black}

      & \cite{WCS+26}
      & 14951.08
      & 179.12
      & 594.24
      & 1577.23
      & 13077.83\\

      & {\bf Ours}
      & \BEST{645.17}
      & \BEST{44.67}
      & \BEST{55.53}
      & \BEST{79.06}
      & \BEST{565.84} \\

    \bottomrule
  \end{tabular}%
  }

  % \caption*{\scriptsize
  % {\cite{CZZ+24} is reported as an RPMT baseline and is excluded
  % from the ranking among ssPMT protocols.}}
\end{table}
%%%%%%%%%%%%%%%%%%%%%%%%PSU%%%%%%%%%%%%%%%%%%%%%%%%%%%%%%%%%%%%%
\paragrap{Comparison of PSU} 
\Cref{tab:psu-comparison} compares our protocol with recent PSU constructions, including RPMT-based protocols that do not prevent DEL~\cite{CZZ+24,PT26} and DEL-free protocols~\cite{JSZG24,TBZ+25,WCS+26}.
% \cite{WCS+26}는 ssPMT 프로토콜 비교에서 분석한 내용처럼 PSU 프로토콜에서도 동일한 경향을 보인다. 
Although \cite{PT26} achieves the lowest raw runtime in most of our experiments, but the construction is not safe on DEL. In particular, its protocol outputs the union to both parties. This two-sided output is inherent in its practical IBLT-peeling procedure: in every round, both parties learn the newly peeled items and use them to update their local IBLTs. Hence, obtaining one-sided output would require more than suppressing one party's final output; it would require additional mechanisms for privately updating the IBLTs. The progressive disclosure of peeled items also means that \cite{PT26} does not prevent during-execution leakage.
Moreover, \cite{PT26} requires $O(\log n)$ rounds to iteratively recover the set difference from an IBLT, whereas our protocol is constant-round. Consequently, \cite{PT26} degrades more sharply as network latency and communication become significant. Our protocol is also less sensitive to the item length: the membership-testing phase can operate on collision-resistant short hashes, whereas the IBLT of \cite{PT26} stores the items themselves and therefore incurs communication linear in their bit length.

\cite{JSZG24} provides protection against DEL, but its heavy communication becomes a severe bottleneck in bandwidth-constrained environments. At \(n=2^{22}\), it communicates \(6186.16\) MB, which is \(8.26\times\) larger than ours; consequently, its runtime at \(100\) Mbps reaches \(586.36\) seconds, compared with \(188.44\) seconds for our protocol. 
Thus, our protocol provides one-sided output and secret-shared intermediate membership results, while offering constant-round execution and substantially lower communication than both \cite{PT26} and \cite{JSZG24}.

\Cref{tab:psu-comparison} compares our protocol with recent PSU constructions. We report the RPMT-based protocols of \cite{CZZ+24,PT26}, which do not prevent DEL, as performance baselines, while \cite{JSZG24,TBZ+25,WCS+26} and ours constitute the directly comparable DEL-free group. 

\cite{PT26} achieves the lowest raw runtime in all LAN experiments and for smaller inputs in some WAN settings. However, it does not provide DEL-free one-sided PSU: its practical IBLT-peeling procedure reveals set items to both parties in each round, resulting in two-sided output and progressive during-execution disclosure. It also requires \(O(\log n)\) rounds to iteratively recover the set difference, whereas our protocol is constant-round. Consequently, our protocol becomes faster than \cite{PT26} in every WAN setting for \(n\ge 2^{20}\). At \(n=2^{22}\), our protocol is \(1.35\times\), \(1.75\times\), and \(2.15\times\) faster at \(1\) Gbps, \(100\) Mbps, and \(10\) Mbps, respectively. Moreover, the IBLT of \cite{PT26} stores the items themselves and therefore incurs communication that grows linearly with the item length, whereas the membership-testing phase of our protocol can operate on collision-resistant short hashes. 

Among the DEL-free protocols, \cite{TBZ+25} achieves the lowest communication for the smaller input sizes, while \cite{JSZG24} achieves the lowest LAN runtime at \(n=2^{16}\) and \(n=2^{18}\). In contrast, for both \(n=2^{20}\) and \(n=2^{22}\), our protocol achieves the lowest communication and runtime in every reported network setting among the DEL-free protocols. At \(n=2^{22}\), our protocol communicates \(748.52\) MB, while the next-best DEL-free protocol, \cite{TBZ+25}, communicates \(1062.79\) MB, or \(1.42\times\) the communication of ours. At \(100\) Mbps, our protocol is \(4.71\times\) faster than \cite{TBZ+25} (\(106.40\) versus \(501.18\) seconds).
The communication bottleneck is more pronounced for \cite{JSZG24} and \cite{WCS+26}: at the same input size, they require \(13.63\times\) and \(20.18\times\) the communication of ours, respectively, and our protocol is \(8.86\times\) and \(14.98\times\) faster than them at \(100\) Mbps. Thus, our construction provides DEL-free one-sided output, constant-round execution, and substantially lower end-to-end costs at large set sizes. The different offline--online cost profiles underlying these total-cost results are examined separately below.


% =========================================================
% 2. PSU comparison table (revised)
% ========================================================
\begin{table}[t]
  \centering
  \caption{Comparison of PSU protocols.
  The best and the runner-up values are marked in red and blue, resp.
  DEL-leaky PSUs (marked with $\ddagger$) are given as baseline, and hence excluded from the colored ranking.\kwon{jszg $2^{16}$ 1gbps 100mbps 재실험중}}
  \label{tab:psu-comparison}

  \footnotesize
  \renewcommand{\arraystretch}{1.15}

  \resizebox{\columnwidth}{!}{%
  \begin{tabular}{ccccccc}
    \toprule
    \multirow{2}{*}{\makecell[c]{Set\\Size}} &
    \multirow{2}{*}{Protocol} &
    \multirow{2}{*}{\makecell{Comm.\\(MB)}} &
    \multicolumn{4}{c}{Runtime (s)} \\
    \cline{4-7}
    & & & 10Gbps & 1Gbps & 100Mbps & 10Mbps \\
    \midrule

    \multirow{6}{*}{$2^{16}$}
      & \cite{CZZ+24}$^\ddagger$
      & 6.78 & 2.47 & 3.66 & 3.82 & 8.26 \\

      & \cite{PT26}$^{\dagger,\ddagger}$
      & 39.09 / 29.74
      & 0.14 & 3.07 & 5.27 & 26.87 \\

    \arrayrulecolor{gray!35}
    \cmidrule{2-7}
    \arrayrulecolor{black}

      & \cite{TBZ+25}
      & \BEST{16.75}
      & 6.64
      & \SECOND{8.79}
      & \SECOND{9.65}
      & \BEST{21.58} \\

      & \cite{JSZG24}
      & 137.42
      & \BEST{1.60}
      & 38.00
      & 29.96
      & 120.00 \\

      & \cite{WCS+26}
      & 234.56
      & 2.61
      & 13.76
      & 28.93
      & 206.47 \\

      & {\bf Ours}
      & \SECOND{38.29}
      & \SECOND{2.50}
      & \BEST{5.96}
      & \BEST{7.85}
      & \SECOND{34.93} \\

    \midrule

    \multirow{6}{*}{$2^{18}$}
      & \cite{CZZ+24}$^\ddagger$
      & 27.07 & 9.83 & 11.61 & 12.97 & 31.01 \\

      & \cite{PT26}$^{\dagger,\ddagger}$
      & 130.07 / 99.87
      & 0.53 & 7.99 & 13.83 & 87.87 \\

    \arrayrulecolor{gray!35}
    \cmidrule{2-7}
    \arrayrulecolor{black}

      & \cite{TBZ+25}
      & \BEST{65.42}
      & 25.59
      & \SECOND{28.63}
      & \SECOND{32.68}
      & \SECOND{112.72} \\

      & \cite{JSZG24}
      & 577.73
      & \BEST{5.47}
      & 37.27
      & 68.42
      & 501.62 \\

      & \cite{WCS+26}
      & 933.85
      & 10.26
      & 40.10
      & 101.81
      & 818.09 \\

      & {\bf Ours}
      & \SECOND{83.92}
      & \SECOND{5.74}
      & \BEST{10.65}
      & \BEST{14.91}
      & \BEST{74.62} \\

    \midrule

    \multirow{6}{*}{$2^{20}$}
      & \cite{CZZ+24}$^\ddagger$
      & 108.23 & 39.31 & 43.41 & 49.99 & 127.09 \\

      & \cite{PT26}$^{\dagger,\ddagger}$
      & 520.18 / 409.04
      & 2.46 & 27.25 & 49.66 & 338.55 \\

    \arrayrulecolor{gray!35}
    \cmidrule{2-7}
    \arrayrulecolor{black}

      & \cite{TBZ+25}
      & \SECOND{262.96}
      & 100.77
      & 108.51
      & \SECOND{125.10}
      & \SECOND{321.20} \\

      & \cite{JSZG24}
      & 2430.47
      & \SECOND{24.55}
      & \SECOND{84.10}
      & 264.74
      & 1346.76 \\

      & \cite{WCS+26}
      & 3731.17
      & 40.90
      & 147.38
      & 393.77
      & 2020.42 \\

      & {\bf Ours}
      & \BEST{227.13}
      & \BEST{16.91}
      & \BEST{25.80}
      & \BEST{35.65}
      & \BEST{198.76} \\

    \midrule

    \multirow{6}{*}{$2^{22}$}
      & \cite{CZZ+24}$^\ddagger$
      & 432.87
      & 159.74
      & 172.78
      & 199.89
      & 514.25 \\

      & \cite{PT26}$^{\dagger,\ddagger}$
      & 2078.78 / 1633.55
      & 10.82
      & 103.97
      & 186.30
      & 1344.22 \\

    \arrayrulecolor{gray!35}
    \cmidrule{2-7}
    \arrayrulecolor{black}

      & \cite{TBZ+25}
      & \SECOND{1062.79}
      & 407.05
      & 433.95
      & \SECOND{501.18}
      & \SECOND{1319.57} \\

      & \cite{JSZG24}
      & 10204.70
      & \SECOND{119.81}
      & \SECOND{330.36}
      & 942.31
      & 5453.89 \\

      & \cite{WCS+26}
      & 15108.02
      & 179.52
      & 599.79
      & 1594.22
      & 13213.29\\

      & {\bf Ours}
      & \BEST{748.52}
      & \BEST{56.59}
      & \BEST{77.24}
      & \BEST{106.40}
      & \BEST{626.12} \\

    \bottomrule
  \end{tabular}%
  }

  \caption*{\scriptsize
  {$\dagger$: \cite{PT26} uses different parameter sets for LAN and WAN
  settings to trade communication for computation.}}
\end{table}
%%%%%%%%%%%%%%%%On/Off%%%%%%%%%%%%%%%%%%%%%%%%%%%%%%

%%%%%%%%%%%%%%%%%%%%%%%%%%
\paragrap{Online latency versus preprocessing cost}
\son{메인 테이블에서 Offline/Online을 안 줄거라서 이게 이렇게 길게, 그리고 앞에 나올 필요는 없다고 생각함. [35]는 계속 total로만 언급해나가다가, 
Online v.s. Offline cost comparison에 가서 "사실 쟤들은 heavy offline을 두고 online을 했다는걸 장점으로 내세움.. 근데 우리도 offline으로 뺄 수 있는 것들 빼면 online 상당히 좋음 + *non-reusable* offline은 실용적으로 크게 의미가 없다고 생각함." }
We report offline, online, and per-execution total costs. 
% Although preprocessing can be completed before the inputs arrive, its correlated state is generally consumed per execution; hence, online cost captures input-dependent latency, whereas total cost reflects the complete resource use. 
% We follow the offline--online separation of each original protocol rather than imposing our own, as preprocessing boundaries and reusability assumptions may differ across implementations.
% \cite{WCS+26} treats permutation-correlation generation and SEQT preprocessing as offline, while hashing, OPRF/OKVS processing, input-dependent equality testing and permutation, and the final OT are online. \cite{JSZG24} instead reports a setup phase comprising base OTs and PEQT preprocessing triples; we therefore count its setup and \emph{w/o setup} times as offline and online, respectively. %\footnote{For \cite{WCS+26}, we completed it  end-to-end wrapper, and for \cite{JSZG24}, we added timing instrumentation for the phase split; neither modification changes the cryptographic submodules.}

% For our protocol, the offline phase includes RLWE key generation and public-key transmission, correlated randomness for the \(NL'\) instances of \(\mathcal{F}_{\mathsf{ss\mbox{-}peqt}}\), and \(L'\) flooded mask ciphertexts. In particular, for each \(\mu\in[L']\), the sender samples \(\vec r_\mu\in\mathbb Z_p^N\) and precomputes a flooded
% \(\mathsf{PlainAdd}(\mathsf{Enc}_{\mathsf{pk}}(\vec 0),\vec r_\mu)\).
% Higher-level protocols additionally preprocess random OTs, permutation-and-share correlations, and comparison-circuit triples; all remaining input-dependent work is online. Protocols reporting only aggregate measurements are listed only under \emph{Total}, since imposing our own split would require protocol-specific assumptions about preprocessing boundaries, reusability, and optimization, thereby undermining a fair comparison. 



%%%%%%%%%%%%%%%%%수정한 버전 작성 파트%%%%%%%%%%%%%
Since \cite{WCS+26} is specifically optimized for low online latency, we separately compare the offline and online costs of the ssPMT protocols for a fair comparison. We preserve the offline--online boundary used by each original work. At a high level, \cite{WCS+26} precomputes permutation and equality-test correlations, \cite{JSZG24} counts its base-OT and PEQT setup as offline, and our protocol precomputes RLWE setup, ss-PEQT correlations, flooded masks, and higher-level OT or P\&S correlations; all remaining input-dependent work is counted as online. 

Online cost measures the latency after the inputs become available, whereas total cost measures the resources consumed per execution. Importantly, input-independent preprocessing is not necessarily reusable: the correlated randomness, secret masks, and OT or P\&S correlations considered here are single-use and must be regenerated, transmitted, and, when multiple executions are prepared in advance, stored separately for every execution. Such preprocessing can reduce input-to-output latency, but it cannot be amortized across executions and therefore still affects provisioning throughput, bandwidth, storage, and total per-execution cost. 

\cite{WCS+26} explicitly optimizes online latency by moving the expensive pSEQT and permutation-correlation components to a single-use offline phase. Applying the same principle to our construction also yields competitive online performance. \kwag{10Mbps에서 online cost 차이는 꽤 큰데 without 10Mbps라고 언급 꼭 해야할까요....} \kwon{wcs 10mbps 재실험중} For \(n=2^{20}\), our online runtime is \(8.16\) seconds at \(1\) Gbps and \(17.35\) seconds at \(100\) Mbps, close to the \(7.92\) and \(16.08\) seconds of \cite{WCS+26}, respectively.
However, the short online phase of \cite{WCS+26} relies on \(3607.72\) MB of offline communication, compared with only \(43.21\) MB in our protocol.Because this preprocessing is single-use, it must be regenerated for every execution. Preparing multiple executions in advance therefore incurs substantial provisioning bandwidth and storage overhead due to the multi-gigabyte correlated state.
At the same bandwidths, our online phase is \(8.92\times\) and \(9.42\times\) faster than that of \cite{JSZG24}. \cite{WCS+26} remains faster online over LAN and at \(10\) Mbps, partly because it communicates less data during the online phase.
 
This online-latency advantage comes with \(3607.72\) MB of single-use offline communication for \cite{WCS+26} at \(n=2^{20}\), \(83.49\times\) that of ours. Overall, \cite{JSZG24} and \cite{WCS+26} require \(10.70\times\) and \(16.43\times\) more total communication, respectively. Consequently, over \(1\) Gbps, \(100\) Mbps, and \(10\) Mbps WANs, our total runtime is \(2.90\times\), \(4.64\times\), and \(6.78\times\) faster than \cite{JSZG24}, and \(5.71\times\), \(11.05\times\), and \(10.17\times\) faster than \cite{WCS+26}. This advantage is particularly relevant in practical WAN settings, where the worldwide median bandwidth is approximately \(100\) Mbps~\cite{InternetSpeed}. Thus, our protocol maintains competitive online latency while substantially reducing non-reusable preprocessing and total per-execution cost.

% =========================================================
% 3. PSU_on/off comparison table (revised)
% ========================================================
\begin{table}[t]
  \centering
  \caption{Comparison of offline and online costs of PSU protocols for $n=2^{20}$.\kwon{jszg $2^{20}$ 10mbps 재실험중}}
  \label{tab:psu-breakdown-220}

  \footnotesize
  \renewcommand{\arraystretch}{1.12}

  \resizebox{\columnwidth}{!}{%
  \begin{tabular}{ccccccc}
    \toprule

    \multirow{2}{*}{Protocol} &
    \multirow{2}{*}{Phase} &
    \multirow{2}{*}{\makecell{Comm.\\(MB)}} &
    \multicolumn{4}{c}{Runtime (s)} \\

    \cline{4-7}

    & & &
    10Gbps & 1Gbps & 100Mbps & 10Mbps \\

    \midrule

    \multirow{3}{*}{\cite{JSZG24}}
      & Offline
      & 252.22
      & 2.95 & 10.49 & 23.74 & 2.45 \\

      & Online
      & 2178.25
      & 21.60 & 73.62 & 240.99 & 1344.31 \\

      & Total
      & 2430.47
      & 24.55 & 84.10 & 264.74 & 1346.76 \\

    \midrule

    \multirow{3}{*}{\cite{WCS+26}}
      & Offline
      & 3607.72
      & 38.16 & 139.46 & 377.70 & 1960.67 \\

      & Online
      & 123.45
      & 2.74 & 7.92 & 16.08 & 59.75 \\

      & Total
      & 3731.17
      & 40.90 & 147.38 & 393.77 & 2020.42 \\

    \midrule

    \multirow{3}{*}{\bf Ours}
      & Offline
      & 43.21
      & 12.74 & 17.64 & 18.29 & 44.79 \\

      & Online
      & 183.91
      & 4.17 & 8.16 & 17.35 & 153.97 \\

      & Total
      & 227.13
      & 16.91 & 25.80 & 35.65 & 198.76 \\

    \bottomrule
  \end{tabular}%
  }
\end{table}
%%%%%%%%%%%%%%%%%%%%%%%%%%%%%%%%%%%%%%%%%%%%


\paragrap{Breakdown of our ssPMT}
\kwag{breakdown 표 업데이트 (준준)}
\Cref{tab:microbenchmark} reports the runtime and communication breakdown of our ssPMT protocol for \(n=2^{20}\) in the LAN setting. As one can easily expect, the largest computational bottleneck is \(\mathsf{RSB\mbox{-}OKVS.HomDecode}\), which is consistent with the cost analysis: each sequenced layer requires \(w+\zeta\) homomorphic operations, and the total work scales as \(O(nw)\). % By contrast, the remaining components are already relatively lightweight, indicating that the practical cost is now concentrated on the core batched homomorphic decoding itself rather than on auxiliary steps such as encryption or decryption.
% \begin{table}[ht]
%     \centering
%     \small 
%     \caption{Breakdown of our RPMT over LAN, for $n=2^{20}$.}
%     \label{tab:microbenchmark}
%     \setlength{\tabcolsep}{6pt} 
    
%     \begin{tabular}{ccc}
%         \toprule
%         \textbf{Stages} & \textbf{Runtime (s)} & \textbf{Comm. (MB)} \\
%         \midrule
%         OPRF & 0.419 & 17.60 \\ \cmidrule(lr){1-3}
%         $\mathsf{RSB\mbox{-}OKVS.Encode}$ & 0.625 & - \\ \cmidrule(lr){1-3}
%         $\mathsf{RLWE.Enc}$ & 0.269 & 55.02 \\ \cmidrule(lr){1-3}
%         $\mathsf{RSB\mbox{-}OKVS.HomDecode}$ & 1.626 & 48.81 \\ \cmidrule(lr){1-3}
%         $\mathsf{RLWE.Dec}$ & 0.032 & - \\ 
%         % $\func{ss\mbox{-}peqt}$ & 11.809 & 45.44 \\\midrule
%         % $\func{ot}$ & 0.271 & 16.64 \\
%         \bottomrule
%     \end{tabular}
% \end{table}
% \begin{table}[ht]
%     \centering
%     \small 
%     \caption{\kwon{online column추가}Breakdown of our ssPMT over LAN, for $n=2^{20}$.}
%     \label{tab:microbenchmark}
%     \setlength{\tabcolsep}{6pt} 
    
%     \begin{tabular}{ccc}
%         \toprule
%         \textbf{Stages} & \textbf{Runtime (s)} & \textbf{Comm. (MB)} \\
%         \midrule
%         Offline & 0.144 & 58.08 \\ \cmidrule(lr){1-3}
%         $\mathsf{RSB\mbox{-}OKVS.Encode}$ & 0.841 & - \\ \cmidrule(lr){1-3}
%         $\mathsf{RLWE.Enc}$ & 0.221 & 48.03 \\ \cmidrule(lr){1-3}
%         $\mathsf{RSB\mbox{-}OKVS.HomDecode}$ & 3.093 & 61.19 \\ \cmidrule(lr){1-3}
%         $\mathsf{RLWE.Dec}$ & 0.085 & - \\ 
%         % $\func{ss\mbox{-}peqt}$ & 11.809 & 45.44 \\\midrule
%         % $\func{ot}$ & 0.271 & 16.64 \\
%         \bottomrule
%     \end{tabular}
% \end{table}
\begin{table}[ht]
    \centering
    \caption{Breakdown of our ssPMT over LAN, for $n=2^{20}$.}
    \label{tab:microbenchmark}

    \footnotesize
    \setlength{\tabcolsep}{3pt}
    \renewcommand{\arraystretch}{1.10}

    \resizebox{\columnwidth}{!}{%
    \begin{tabular}{c|c|cc}
        \toprule

        \multicolumn{2}{c|}{\textbf{Stages}} &
        \textbf{Runtime (s)} &
        \textbf{Comm. (MB)} \\

        \midrule

        \multicolumn{2}{c|}{Offline}
        & 9.330
        & 31.88 \\

        \midrule

        \multirow{5}{*}{Online}
          & $\mathsf{RSB\mbox{-}OKVS.Encode}$
          & 0.748
          & - \\

          & $\mathsf{RLWE.Enc}$
          & 0.223
          & 48.03 \\

          & $\mathsf{RSB\mbox{-}OKVS.HomDecode}$
          & 3.027
          & 61.18 \\

          & $\mathsf{RLWE.Dec}$
          & 0.084
          & - \\

          & $\mathcal{F}_{\mathsf{ss\mbox{-}peqt}}$
          & 0.039
          & 58.09 \\

        \midrule

        \multicolumn{2}{c|}{\textbf{Total}}
        & 13.451
        & 199.17 \\

        \bottomrule
    \end{tabular}%
    }
\end{table}

%%%%%%%%%%%%%%%%%%%%%%%%%%%%%%%%%%%%%

% \begin{table*}[t]
%   \centering
%   \caption{Comparison of ssPMT protocols.
%   For \cite{CZZ+24}, we report its RPMT protocol.
%   All timing reports include setup timing.}
%   \label{tab:sspmt}

%   \small
%   \setlength{\tabcolsep}{0.5pt}
%   \renewcommand{\arraystretch}{1.00}

%   \begin{tabular}{@{}c|c|ccccc|ccccc|ccccc@{}}
%     \toprule

%     \multirow{3}{*}{\makecell{Set\\Size}} &
%     \multirow{3}{*}{Protocol} &
%     \multicolumn{5}{c|}{Offline} &
%     \multicolumn{5}{c|}{Online} &
%     \multicolumn{5}{c}{Total} \\

%     \cmidrule(lr){3-7}
%     \cmidrule(lr){8-12}
%     \cmidrule(lr){13-17}

%     & &
%     \multirow{2}{*}{\makecell{Comm.\\(MB)}} &
%     \multicolumn{4}{c|}{Runtime (s)} &
%     \multirow{2}{*}{\makecell{Comm.\\(MB)}} &
%     \multicolumn{4}{c|}{Runtime (s)} &
%     \multirow{2}{*}{\makecell{Comm.\\(MB)}} &
%     \multicolumn{4}{c}{Runtime (s)} \\

%     \cmidrule(lr){4-7}
%     \cmidrule(lr){9-12}
%     \cmidrule(lr){14-17}

%     & & &
%     10Gbps & 1Gbps & 100Mbps & 10Mbps &
%     & 10Gbps & 1Gbps & 100Mbps & 10Mbps &
%     & 10Gbps & 1Gbps & 100Mbps & 10Mbps \\

%     \midrule

%     \multirow{3}{*}{$2^{16}$}
%       & \cite{CZZ+24}
%       & - & - & - & - & -
%       & - & - & - & - & -
%       & 4.67 & 2.45 & 3.31 & 3.46 & - \\

%     \arrayrulecolor{gray!35}
%     \cmidrule(lr){2-17}
%     \arrayrulecolor{black}

%       & \cite{WCS+26}
%       & 226.19 & 2.41 & 11.46 & 26.29 & 198.07
%       & 8.24 & 0.16 & 2.01 & 2.35 & 8.13
%       & 234.44 & 2.57 & 13.47 & 28.64 & 206.21 \\

%       & \textbf{Ours}
%       & 5.96 & 1.71 & 8.09 & 7.38 & 12.92
%       & 29.25 & 0.23 & 1.53 & 3.47 & 25.25
%       & 35.21 & 1.94 & 9.62 & 10.85 & 38.17 \\

%     \midrule

%     \multirow{3}{*}{$2^{18}$}
%       & \cite{CZZ+24}
%       & - & - & - & - & -
%       & - & - & - & - & -
%       & 18.66 & 9.80 & 11.12 & 11.95 & - \\

%     \arrayrulecolor{gray!35}
%     \cmidrule(lr){2-17}
%     \arrayrulecolor{black}

%       & \cite{WCS+26}
%       & 902.49 & 9.57 & 36.82 & 96.44 & 789.85
%       & 31.22 & 0.60 & 2.94 & 5.00 & 27.79
%       & 933.71 & 10.17 & 39.76 & 101.44 & 817.64 \\

%       & \textbf{Ours}
%       & 12.30 & 3.94 & 16.86 & 17.74 & 27.30
%       & 63.15 & 0.92 & 2.63 & 6.81 & 54.47
%       & 75.45 & 4.86 & 19.48 & 24.55 & 81.77 \\

%     \midrule

%     \multirow{3}{*}{$2^{20}$}
%       & \cite{CZZ+24}
%       & - & - & - & - & -
%       & - & - & - & - & -
%       & 74.66 & 39.20 & 42.14 & 46.59 & - \\

%     \arrayrulecolor{gray!35}
%     \cmidrule(lr){2-17}
%     \arrayrulecolor{black}

%       & \cite{WCS+26}
%       & 3607.72 & 38.16 & 139.46 & 377.70 & 1960.67
%       & 123.29 & 2.36 & 7.22 & 15.34 & 58.90
%       & 3731.01 & 40.52 & 146.69 & 393.04 & 2019.57 \\

%       & \textbf{Ours}
%       & 31.88 & 9.56 & 44.29 & 27.42 & 71.51
%       & 167.30 & 4.11 & 7.07 & 16.84 & 144.60
%       & 199.17 & 13.67 & 51.36 & 44.26 & 216.11 \\

%     \midrule

%     \multirow{3}{*}{$2^{22}$}
%       & \cite{CZZ+24}
%       & - & - & - & - & -
%       & - & - & - & - & -
%       & 298.63 & 158.97 & 168.52 & 187.18 & - \\

%     \arrayrulecolor{gray!35}
%     \cmidrule(lr){2-17}
%     \arrayrulecolor{black}

%       & \cite{WCS+26}
%       & & & & &
%       & & & & &
%       & & & & & \\

%       & \textbf{Ours}
%       & 99.11 & 28.28 & 136.62 & 81.99 & 178.38
%       & 546.06 & 15.30 & 23.86 & 54.44 & 474.32
%       & 645.17 & 43.58 & 160.48 & 136.43 & 652.70 \\

%     \bottomrule
%   \end{tabular}
% \end{table*}



%%%%%%%%%%%%%%%%%%%%%%%%%%%%%%%%%%%%%%%%%%%%%%%%%%%%%%

  % \caption{Comparison of PSU protocols.
  % All timing reports include setup timing. \kwon{WCS 10Mbps $2^{20}$ 이상치 발견. 재실험필요} \kwag{ [17]은 commu만 on,offline 작성안되어 있는데 이유가 뭐지? (준준)}}



% \begin{table*}[t]
%   \centering
%   \caption{Comparison of PSU protocols.
%   All timing reports include setup timing.}
%   \label{tab:psu}

%   \small
%   \setlength{\tabcolsep}{0.45pt}
%   \renewcommand{\arraystretch}{1.00}

%   \begin{tabular}{@{}c|c|ccccc|ccccc|ccccc@{}}
%     \toprule

%     \multirow{3}{*}{\makecell{Set\\Size}} &
%     \multirow{3}{*}{Protocol} &
%     \multicolumn{5}{c|}{Offline} &
%     \multicolumn{5}{c|}{Online} &
%     \multicolumn{5}{c}{Total} \\

%     \cmidrule(lr){3-7}
%     \cmidrule(lr){8-12}
%     \cmidrule(lr){13-17}

%     & &
%     \multirow{2}{*}{\makecell{Comm.\\(MB)}} &
%     \multicolumn{4}{c|}{Runtime (s)} &
%     \multirow{2}{*}{\makecell{Comm.\\(MB)}} &
%     \multicolumn{4}{c|}{Runtime (s)} &
%     \multirow{2}{*}{\makecell{Comm.\\(MB)}} &
%     \multicolumn{4}{c}{Runtime (s)} \\

%     \cmidrule(lr){4-7}
%     \cmidrule(lr){9-12}
%     \cmidrule(lr){14-17}

%     & & &
%     10Gbps & 1Gbps & 100Mbps & 10Mbps &
%     & 10Gbps & 1Gbps & 100Mbps & 10Mbps &
%     & 10Gbps & 1Gbps & 100Mbps & 10Mbps \\

%     \midrule

%     \multirow{6}{*}{$2^{16}$}
%       & \cite{CZZ+24}
%       & - & - & - & - & -
%       & - & - & - & - & -
%       & 6.78 & 2.47 & 3.66 & 3.82 & 8.26 \\

%       & \cite{PT26}
%       & - & - & - & - & -
%       & - & - & - & - & -
%       & 39.09/29.74$^\dagger$
%       & 0.14 & 3.07 & 5.27 & 26.87 \\

%     \arrayrulecolor{gray!35}
%     \cmidrule(lr){2-17}
%     \arrayrulecolor{black}

%       & \cite{TBZ+25}
%       & - & - & - & - & -
%       & - & - & - & - & -
%       & 16.75 & 6.64 & 8.79 & 9.65 & 21.58 \\

%       & \cite{JSZG24}
%       & - & 0.22 & 1.97 & 1.97 & 1.97
%       & - & 0.72 & 37.54 & 42.11 & 82.57
%       & 97.46 & 0.94 & 39.51 & 44.08 & 84.54 \\

%       & \cite{WCS+26}
%       & 226.19 & 2.41 & 11.46 & 26.29 & 198.07
%       & 8.37 & 0.20 & 2.29 & 2.65 & 8.39
%       & 234.56 & 2.61 & 13.76 & 28.93 & 206.47 \\

%       & \textbf{Ours}
%       & 7.95 & 2.46 & 10.01 & 10.43 & 16.61
%       & 30.35 & 0.23 & 1.20 & 3.44 & 26.13
%       & 38.29 & 2.70 & 11.20 & 13.87 & 42.74 \\

%     \midrule

%     \multirow{6}{*}{$2^{18}$}
%       & \cite{CZZ+24}
%       & - & - & - & - & -
%       & - & - & - & - & -
%       & 27.07 & 9.83 & 11.61 & 12.97 & 31.01 \\

%       & \cite{PT26}
%       & - & - & - & - & -
%       & - & - & - & - & -
%       & 130.07/99.87$^\dagger$
%       & 0.53 & 7.99 & 13.83 & 87.87 \\

%     \arrayrulecolor{gray!35}
%     \cmidrule(lr){2-17}
%     \arrayrulecolor{black}

%       & \cite{TBZ+25}
%       & - & - & - & - & -
%       & - & - & - & - & -
%       & 65.42 & 25.59 & 28.63 & 32.68 & 112.72 \\

%       & \cite{JSZG24}
%       & - & 0.25 & 1.78 & 1.78 & 1.96
%       & - & 3.50 & 45.07 & 53.80 & 335.24
%       & 387.71 & 3.75 & 46.84 & 55.58 & 337.20 \\

%       & \cite{WCS+26}
%       & 902.49 & 9.57 & 36.82 & 96.44 & 789.85
%       & 31.36 & 0.69 & 3.28 & 5.36 & 28.23
%       & 933.85 & 10.26 & 40.10 & 101.81 & 818.09 \\

%       & \textbf{Ours}
%       & 16.55 & 5.15 & 19.80 & 18.23 & 33.61
%       & 67.37 & 0.94 & 2.51 & 7.04 & 58.09
%       & 83.92 & 6.09 & 22.31 & 25.27 & 91.70 \\

%     \midrule

%     \multirow{6}{*}{$2^{20}$}
%       & \cite{CZZ+24}
%       & - & - & - & - & -
%       & - & - & - & - & -
%       & 108.23 & 39.31 & 43.41 & 49.99 & 127.09 \\

%       & \cite{PT26}
%       & - & - & - & - & -
%       & - & - & - & - & -
%       & 520.18/409.04$^\dagger$
%       & 2.46 & 27.25 & 49.66 & 338.55 \\

%     \arrayrulecolor{gray!35}
%     \cmidrule(lr){2-17}
%     \arrayrulecolor{black}

%       & \cite{TBZ+25}
%       & - & - & - & - & -
%       & - & - & - & - & -
%       & 262.96 & 100.77 & 108.51 & 125.10 & 321.20 \\

%       & \cite{JSZG24}
%       & - & 0.42 & 1.89 & 2.05 & 2.45
%       & - & 15.23 & 72.80 & 163.41 & 1344.31
%       & 1547.70 & 15.64 & 74.69 & 165.46 & 1346.76 \\

%       & \cite{WCS+26}
%       & 3607.72 & 38.16 & 139.46 & 377.70 & 1960.67
%       & 123.45 & 2.74 & 7.92 & 16.08 & 59.75
%       & 3731.17 & 40.90 & 147.38 & 393.77 & 2020.42 \\

%       & \textbf{Ours}
%       & 43.21 & 13.05 & 50.80 & 54.06 & 86.33
%       & 183.91 & 4.14 & 7.73 & 18.46 & 160.19
%       & 227.13 & 17.19 & 58.53 & 72.52 & 246.52 \\

%     \midrule

%     \multirow{6}{*}{$2^{22}$}
%       & \cite{CZZ+24}
%       & - & - & - & - & -
%       & - & - & - & - & -
%       & 432.87 & 159.74 & 172.78 & 199.89 & 514.25 \\

%       & \cite{PT26}
%       & - & - & - & - & -
%       & - & - & - & - & -
%       & 2078.78/1633.55$^\dagger$
%       & 10.82 & 103.97 & 186.30 & 1344.22 \\

%     \arrayrulecolor{gray!35}
%     \cmidrule(lr){2-17}
%     \arrayrulecolor{black}

%       & \cite{TBZ+25}
%       & - & - & - & - & -
%       & - & - & - & - & -
%       & 1062.79 & 407.05 & 433.95 & 501.18 & 1319.57 \\

%       & \cite{JSZG24}
%       & - & 1.26 & 2.69 & 3.06 & 4.46
%       & - & 62.80 & 228.97 & 583.30 & 5449.43
%       & 6186.16 & 64.05 & 231.66 & 586.36 & 5453.89 \\

%       & \cite{WCS+26}
%       & & & & &
%       & & & & &
%       & & & & & \\

%       & \textbf{Ours}
%       & 136.47 & 40.82 & 155.76 & 127.27 & 268.18
%       & 612.05 & 15.41 & 25.91 & 61.17 & 530.44
%       & 748.52 & 56.22 & 181.68 & 188.44 & 798.63 \\

%     \bottomrule
%   \end{tabular}

%   \caption*{\small
%   {$\dagger$: \cite{PT26} uses different parameter sets for LAN and WAN
%   settings to trade communication for computation.}}
% \end{table*}



%%%%%%%%%%%%%%%%%%%%%%%%%%%%%%%%%%%%%%%%%%%%%%%%
\section{Conclusion}
We presented an efficient ssPMT protocol based on RLWE-batched homomorphic OKVS decoding. Our main technical contribution is RSB-OKVS, obtained through a structure-preserving column permutation of random-band OKVS, together with a span-constrained sequencing algorithm that arranges colliding query rows into SIMD-friendly layers. By padding the input-dependent number of layers to a public budget $L'$ and applying masking, rerandomization, and noise flooding before secret-shared equality testing, our protocol produces a fixed-length collection of XOR-shared membership indicators in a protocol-defined layout. These shares can be directly used to construct several private set operations, including PSI-Cardinality, PSI-Threshold, PSI-Sum, and PSU.
Furthermore, regarding the computation of the inner product over the intersection—where both parties possess associated payloads—several dedicated solutions exist~\cite{PSTY19,LPR+21}. 
% Extending our approach to efficiently support this functionality remains an open problem.

\newpage
%%%%%%%%%%%%%%%%%%%%%%%%%%%%%%%%%%%%%%%%%%%%%%%%

%%%%%%%%%%%%%%%%%%%%%%%%%%%%%%%%%%%%%%%%%%%%%%%
% \section*{Ethical Considerations}
% This work does not involve human participants, user studies, or the collection of real-world personal/operational data.
% All datasets used in this paper are generated solely for our experiments via a fully computational process, and do not incorporate any pre-existing or real-world datasets.
% We therefore do not raise concerns regarding informed consent, personal-data privacy, or confidentiality.
% Our experiments are conducted offline in a controlled environment and do not interact with third-party production systems.
% The released materials contain no sensitive information; consequently, the residual risk of harm to stakeholders is minimal.

% \section*{Open Science}
% We follow with Open Science principles to promote the transparency and reproducibility of our research. 
% To support double-blind review, we provide our research artifacts in an anonymous open-source GitHub repository at \url{https://anonymous.4open.science/r/RLWE-OKVS-522E}
% This repository provides the complete source code, scripts for data generation, and tools for experimental analysis, enabling full reproduction of the results reported in this paper. Detailed setup is in the repository README. Once the paper is accepted, we will publicly release our research artifacts in an open-source GitHub repository. By releasing our work, we aim to enable follow-up research and advance a more open and collaborative scientific environment.

% \section*{Acknowledgements}
% This paper was edited for grammar and translation using ChatGPT.

% %%%%%%%%%%%%%%%%%%%%%%%%%5

% ---- Bibliography ----
%
% BibTeX users should specify bibliography style 'splncs04'.
% References will then be sorted and formatted in the correct style.
%

\bibliographystyle{plainurl}
\bibliography{mybibliography}

%%
%% If your work has an appendix, this is the place to put it.
\appendix


\section*{Ethical Considerations}
\textbf{Papers in-scope for USENIX Security frequently, if not always, have both ethical considerations and implications that should be addressed prior to conducting research. We strongly encourage that you use this appendix, up to one page, to explain the ethical considerations of your work.}


\section*{Open Science}
\textbf{Within up to one page, this appendix must list all artifacts necessary to evaluate the contribution of the paper and make clear how the review committees can access each artifact. See the Call for Papers for details.}


\section{Ideal Functionalities}\label{sec:Funtionality}
The ideal functionalities for oblivious transfer, oblivious PRF, and Secret-Shared Private Equality Test are as follows.
%%%%%%%%%%%%%%%%%%%%%%%%%%%%
% \begin{figure}[htbp]
% \centering
% \fbox{%
% \begin{minipage}{0.95\linewidth}
% \textbf{Parameters:} Sender $\sender$, receiver $\receiver$, set size $n$, and an encryption scheme 
% $\mathcal{E}=({\sf Setup},{\sf KeyGen},{\sf Enc},{\sf Dec})$.\\
% \textbf{Functionality:}
% \begin{itemize}
%     \item Wait for input $k$ and $s$ from the receiver $\receiver$.
%     \item Wait for input $\{s_1^*,\cdots,s_n^*\}\subseteq \{0,1\}^*$ from the sender $\sender$.
%     \item For each $i\in[n]$:
%     \begin{itemize}
%         \item Compute $s_i' \gets {\sf Dec}(k,s_i^*)$.
%         \item If $s_i'=s$, set $b_i=1$; otherwise set $b_i=0$.
%     \end{itemize}
%     \item Give output $\vec b=(b_1,\cdots,b_n)\in\{0,1\}^n$ to the receiver $\receiver$.
% \end{itemize}
% \end{minipage}
% }
% \caption{\blue{Vector Oblivious Decryption-then-Matching Functionality $\mathcal{F}_{\sf vodm}$.}}
% \label{fig:vodm_functionality}
% \end{figure}

% {\color{blue}The RPMT construction based on VODM~\cite{ZCL+23} is depicted in \Cref{fig:zcl_pke_framework}. This is secure for the following reasons:
% \begin{itemize}
%     \item For the case when $y_j \in X$, the OKVS-decoded ciphertext $c_j^*$ would decrypt to {\sf IS}. However, if one sends $c_j^*$ as it is, $\receiver$ can find the exact same $c_j^*$ and learn which $x_i$ is contained in the other party's set $Y$, which is problematic. The re-randomization step converts $c_j^*$ into another ciphertext while maintaining the inner message, resolving this issue. Concretely, it is implemented by generating a ciphertext of $0$ (an encryption of $0$) by itself and adding it.
    
%     \item For the case when $y_j \notin X$, observe that the randomized ciphertext $c_j^*$ still contains information about $y_j$, since the inner message itself remains unchanged even after re-randomization. Thus, rigorously speaking, their construction is not a secure realization of RPMT.\footnote{The authors were also aware of this, and call this construction under the leaky-VODM and leaky-RPMT.} However, since the real goal of \cite{ZCL+23} was private set union that aims to reveal $Y-X$ to the receiver, such partial information about the element $y_j \notin X$ is not a leakage at all, and hence this leaky construction was fine. 
% \end{itemize}
% }

%%%%%%%%%%%%%%%%%%%%%%%%%%%%%%%%%
\begin{figure}[htbp]
\centering
\fbox{%
\begin{minipage}{0.95\linewidth}
\textbf{Parameters:} Sender $\mathcal{S}$, Receiver $\mathcal{R}$, message length $\ell_m$ \\
\textbf{Input:} $b \in \{0,1\}$ of $\receiver$, and $m_0, m_1 \in\{0,1\}^{\ell_m}$ of $\sender$\\
\textbf{Functionality:} Output $m_b$ to $\mathcal{R}$.
\end{minipage}%
}
\caption{(1-out-of-2) Oblivious transfer functionality $\mathcal{F}_{\mathsf{ot}}$.}\label{fig:ot}
\end{figure}
%%%%%%%%%%%%%%%%%%%%%%%%%%%%%%%%%%%%%%%%%%%5

% %%%%%%%%%%%%%%%%%%%%%%%%%%%5
\begin{figure}[h]
\centering
\fbox{%
\begin{minipage}{0.95\linewidth}
\textbf{Parameters:} Sender $\mathcal{S}$, Receiver $\mathcal{R}$, message length $\ell_m$ \\
\textbf{Input:} $\vec x = (x_i)$ from $\sender$, $\vec y = (y_i)$ from $\receiver$\\
\textbf{Functionality:} Define a bit-string $\vec b = \left(\vec 1(x_i = y_i)\right)$. Sample a random bit-string $\vec s$, and output $\vec s$ to $\sender$, and $\vec s\oplus \vec b$ to $\receiver$.
\end{minipage}%
}
\caption{Ideal functionality $\mathcal{F}_{\mathsf{ss\mbox{-}peqt}}$ for secret-shared private equality test.}\label{fig:ss-peqt}
\end{figure}

%%%%%%%%%%%%%%%%%%%%%%%%%%%%%%%%%%%%%%%%%%
\section{Generating Homomorphic Decoding Matrix: Example}\label{sec:example}
To illustrate the process in \Cref{sec:rlwe-okvs}, we provide the examples as follows. For clarity, we assume small parameter $n=8, m=20, w=3, N=4$, and $\zeta=2$.

According to \Cref{eq:N_perm_mat}, random band (RB) matrix \Cref{fig:RBa} is converted to random spaced band (RSB) matrix \Cref{fig:PermutedRSB}. 
For example, the second row starting point $i=9$ in RB matrix corresponds to $(i_s,i_b)=(2,4)$; thus $\pi_N(9)=(4-1)\cdot 4+2=14$ in RSB matrix.

%%%%%%%%%%%%%%%%%%Figure%%%%%%%%%%%%%%%%%%%%%%%%%

\begin{figure}[htbp]
    \centering
    %%%%%
    \begin{subfigure}[b]{0.23\textwidth}
        \centering
        \includegraphics[width=\textwidth]{images/1.png}
        \caption{Random band matrix}\label{fig:RBa}
    \end{subfigure}
    \hfill
    %%%%
    \begin{subfigure}[b]{0.23\textwidth}
        \centering
        \includegraphics[width=\textwidth]{images/2.png}
        \caption{Random spaced band matrix}\label{fig:PermutedRSB}
    \end{subfigure}
    \caption{Structural comparison between a RB matrix and RSB matrix under the column permutation $\pi_N$.}
    \label{fig:PermutatedRSBMatrix}
\end{figure}
%%%%%%%%%%%%%%%%%%%%%%%%%%%%%%%%%

Second, the compact sequencing with clustered rows proceeds as illustrated in \Cref{fig:Sequencing}.
The first row (starting position $a_1=1$, block index $b_1=1$) is placed in the first $N\times m$ layer matrix $\mat{\bar{M}}_1$, but when processing the second row ($a_2=14$, $b_2=4$), the block distance exceeds the span limit ($|b_2-b_1|>\zeta$), so we create a new layer $\mat{\bar{M}}_2$.
Subsequently, the third row ($a_3=1$, $b_3=1$) is also allocated to a newly created third layer $\mat{\bar{M}}_3$. 
This is because it collides with the first row in $\mat{\bar{M}}_1$, and cannot be placed in $\mat{\bar{M}}_2$ due to the span limit ($|b_3-b_2| > \zeta$).

%%%%%%%%%%% our sequencing Fig %%%%%%%%%%%%%%%%%%%%%
\begin{figure}[htbp]
    \centering
    \begin{subfigure}[h]{0.48\textwidth}
     \vspace{0pt}
        \centering
        \includegraphics[width=\textwidth]{images/our seq1.JPG}
        \vspace{5pt}
        % \caption{The second row starts a new layer, as it fits in no existing layer.}
    \end{subfigure}\hfill
    \begin{subfigure}[h]{0.48\textwidth}
     \vspace{0pt}
        \centering
        \includegraphics[width=0.61\textwidth]{images/our seq2.JPG}
        \includegraphics[width=0.37\textwidth]{images/our seq3.JPG}
        \vspace{5pt}
    \end{subfigure}\hfill
    % \begin{subfigure}[h]{0.44\textwidth}
    %  \vspace{0pt}
    %     % \centering
    %      % \includegraphics[width=0.35\textwidth]{images/our seq1.JPG}        
    %     \hfill\includegraphics[width=0.56\textwidth]{images/our seq3.JPG}
    %     % \caption{The third row also requires a new layer, because there is no layer where it can be placed. The final sequencing result is given as the rightmost figure.}
    % \end{subfigure}

    \caption{Our compact sequencing process.}
    \label{fig:Sequencing}
\end{figure}

%%%%%%%%%%%%%%%%%%%%%%%%%%
\section{{\red {Formal Proof for \Cref{thm:semi-honest-rl-sspmt}}}}\label{sec:security_proof}
\kwag{수정 필요}
\paragrap{Formal definition of semi-honest security}
Consider a two-party protocol~$\Pi$ that computes an ideal functionality $f(X, Y)$ where party $P_1$ has input $X$ and $P_2$ has input $Y$.
For each party $P_i$, let ${\sf view}_i(X, Y)$ denote the view of $P_i$ during an execution of $\Pi$ on input $X, Y$.
Specifically, it consists of the input $X$ (or $Y$), the messages sent or received, the random tape sampled during the protocol execution, and the output of the protocol~$\Pi$.

\begin{definition}
We say that a (two-party) protocol $\Pi$ for $P_1$ and $P_2$ computes $f$ securely against semi-honest adversaries if there exist simulators ${\sf Sim}_1$ and ${\sf Sim}_2$ such that for any inputs $X, Y$ of $P_1$ and $P_2$, respectively, 
\begin{align*}
    {\sf Sim}_1(X, f(X, Y)) \cong_c {\sf view}_1(X, Y) \\
{\sf Sim}_2(Y, f(X, Y)) \cong_c {\sf view}_2(X, Y) 
\end{align*}
where $\cong_c$ represents computational indistinguishability.
\end{definition}

\paragrap{Simulation for the corrupted sender}
In the ideal execution of $\mathcal{F}_{\mathsf{RPMT}}$, $\sender$ has no output, and hence the simulator given only $\sender$'s input $Y$. It behaves identically to the real protocol, except send a dummy ciphertext of a random vector $\tilde{\vec p}$, instead of the OKVS encoding of $X'$. 
This simulation is clearly indistinguishable from the real view, thanks to the IND-CPA security of RLWE.

\paragrap{Simulation for the corrupted receiver}
In the ideal execution of $\mathcal{F}_{\mathsf{RPMT}}$, $\receiver$ is given the membership bit vectors $\vec b = (b_i)_{i\in[n_y]}$ where $b_i = \vec 1(y_i \in X)$, and hence the simulator is given $\receiver$'s input $X$ and the ideal output vector $\vec b$. The simulator is defined to run the protocol with the following fake OPRF set $\widetilde{Y}'$:
Compute $F_k(X)$ by itself, randomly samples a $|X\cap Y|$-sized subset among $F_k(X)$, fill the set with a random element until reaching to size $n_y$, and finally reorder the fake set $\widetilde{Y}'$ so that $y_i' \in F_k(X)$ iff $b_i = 1$.

This simulation is indistinguishable from the real protocol's view. From $\receiver$'s point of view, $\sender$-side protocol after OPRF depends on $Y$ only through the randomized labels in $Y'$.
For indices $i$ with $b_i=1$, all that matters for correctness is that the corresponding randomized value belongs to $X'$, since every key in the OKVS table is associated with the same indicator value $\mathsf{IS}$.
For indices $i$ with $b_i=0$, the randomized values are pseudorandom labels outside $X'$, and hence are computationally indistinguishable from fresh random labels.
Therefore, conditioned on $X'$ and the ideal output vector $\vec b$, the real randomized sender set $Y'$ is computationally indistinguishable from the fake set $\widetilde{Y}'$.$\square$

%%%%%%%%%%%%%%%%%%%%%%%%%%%%%%%%%%%%%%%%
\section{Proof for Lemma~\ref{lemma:Lprime}} \label{app:Lprime}

Write each starting point as
\[
  a_i=(\beta_i-1)N+\tau_i,
  \qquad
  \beta_i\in[b],\quad \tau_i\in[N],
\]
where $\beta_i$ and $\tau_i$ denote its block and bin, respectively. For
a block interval $J$, define
\[
  n_\tau(J) := \bigl|\{i:\tau_i=\tau,\ \beta_i\in J\}\bigr|,
  \qquad
  M(J):=\max_{\tau\in[N]}n_\tau(J),
\]
and let $\operatorname{ext}(J) := [\min J-\zeta+1,\max J]$, so that
$|\operatorname{ext}(J)|=g+\zeta-1$ when $|J|=g$. Call a family
$J_1,\ldots,J_K$ of block intervals \emph{admissible} if the extensions
$\operatorname{ext}(J_k)$ are pairwise disjoint, and define
\[
  D(u) \;=\; \max\Big\{ \textstyle\sum_{k=1}^{K} M(J_k) \;:\;
     J_1,\ldots,J_K \text{ admissible},\ \max_k \max J_k \le u \Big\},
\]
with the empty family (value $0$) allowed.

We use three properties of \Cref{alg:sequencing}: rows are inserted in
nondecreasing block order, into the admissible layer of smallest index.
\begin{itemize}
  \item[(P1)] A layer's \emph{anchor}, the block of the row that opened
        it, is the least block it ever contains; anchors are created in
        nondecreasing order, and a layer with anchor $s$ accepts exactly
        the blocks $[s,\, s+\zeta-1]$.
  \item[(P2)] A row $(\tau,\beta)$ scans the existing layers with anchor
        in $[\beta-\zeta+1,\beta]$ in order of anchor and joins the
        first one whose bin $\tau$ is free. In particular, if it settles
        at anchor $s$, every layer existing at that moment with anchor
        in $[\beta-\zeta+1,\, s-1]$ was $\tau$-occupied at that moment.
  \item[(P3)] An occupied (layer, bin) slot never vacates, and a layer
        holds at most one row per bin.
\end{itemize}

\begin{proposition}\label{prop:det-Lprime}
On every input, the number of layers opened by \Cref{alg:sequencing} is
at most $D(b)$.
\end{proposition}

The key is a structural property of the moments at which a new layer is
opened. Fix such a moment: the algorithm, holding the layer multiset
$\Lambda$, fails to place a row of bin $\tau$ at block $y$. Write
$a(\ell)$ for the anchor of a layer $\ell$.

\begin{lemma}[Blocking structure]\label{lem:blocking-Lprime}
There is an $X \le y-\zeta+1$ such that (i) the layers of $\Lambda$ with
$a(\ell)\ge X$ all have $a(\ell)\in[X,y]$, and each of them is
$\tau$-occupied by a row with block in $[X+\zeta-1,\,y]$; (ii) with
$K=|\{\ell\in\Lambda : a(\ell)\ge X\}|$ and counting the failing row
itself, $M([X+\zeta-1,\,y]) \ge K+1$.
\end{lemma}

\begin{proof}
We build a closure. Maintain a collection of \emph{ranges} of anchors,
initially the single range $[y-\zeta+1,\,y]$; whenever a layer
$\ell\in\Lambda$ has its anchor inside a known range, absorb it: we
argue below that $\ell$ is $\tau$-occupied, by a unique row (its
\emph{witness}, unique by (P3)) with block $\beta_\ell$; add the range
$[\beta_\ell-\zeta+1,\;s_\ell]$, where $s_\ell = a(\ell)$. Iterate to a
fixpoint, and let $X$ be the least left endpoint of any range.

First, every absorbed layer is $\tau$-occupied. We induct over the
absorption order. A layer absorbed by the initial range was admissible
to the failing row, which failed. Otherwise $\ell$'s anchor lies in
$[\beta_p-\zeta+1,\, s_p]$ for an earlier witness $p$. If
$a(\ell) < s_p$: by (P1), $\ell$ was opened while block
$a(\ell) < s_p \le \beta_p$ was current, hence before $p$ was
processed; at $p$'s turn it existed, was admissible
($a(\ell)\ge\beta_p-\zeta+1$), and preceded $p$'s layer in the scan, so
by (P2) it was $\tau$-occupied then, hence still, by (P3). If
$a(\ell)=s_p$: either $\ell$ is $p$'s own layer, occupied by $p$, or we
trace how $p$'s layer was absorbed. If through a range containing $s_p$
strictly below its right endpoint, the previous case applies to $\ell$
verbatim; if through the initial range, then
$a(\ell)\in[y-\zeta+1,\,y]$ and the failure applies; if through another
witness's endpoint $s_q=s_p$, we ascend to $q$. The ascent decreases
the absorption index, so it terminates in a resolved case.

Second, the ranges cover exactly $[X,y]$. A new range
$[\beta_\ell-\zeta+1,\, s_\ell]$ contains $s_\ell$, since a row's block
is within $\zeta-1$ of its layer's anchor, and $s_\ell$ lies in the
range that absorbed $\ell$; so the union stays a single interval, and
its right end stays $y$ because no anchor exceeds the current block.

Consequently every layer with $a(\ell)\ge X$ has $a(\ell)\in[X,y]$, is
covered by a range, and was absorbed, which gives (i): each witness has
$\beta_\ell-\zeta+1\ge X$ by the minimality of $X$, i.e.\
$\beta_\ell\ge X+\zeta-1$. For (ii), the witnesses are pairwise-distinct
rows of bin $\tau$ with blocks in $[X+\zeta-1,\,y]$; together with the
failing row (block $y \ge X+\zeta-1$) they number $K+1$.
\end{proof}

\begin{proof}[Proof of \Cref{prop:det-Lprime}]
Write $G(u)$ for the number of layers with anchor $\le u$ opened so
far. We show by induction over the processed rows that $G(u)\le D(u)$
for every $u$, where $G$, $D$, $M$ all refer to the current prefix of
rows. Initially both sides vanish. Placing a row into an existing layer
changes no $G(u)$, and never decreases $D(u)$ since an added row only
increases the counts $M(J)$. Suppose a row of bin $\tau$ at block $y$
opens a new layer, and let $X$, $K$ be as in
\Cref{lem:blocking-Lprime}. For $u<y$ nothing changes, as the new
anchor is $y$. For $u\ge y$, take an admissible family with right
endpoints $\le X-1$ certifying $G(X-1)$ by induction (empty if
$X\le1$), and adjoin $J_0=[X+\zeta-1,\,y]$: its extension
$\operatorname{ext}(J_0)=[X,y]$ is disjoint from the previous
extensions $\subseteq(-\infty,X-1]$, and $M(J_0)\ge K+1$ by the lemma,
counting the newly placed row. The family's value is now at least
$G(X-1)+K+1$, which by \Cref{lem:blocking-Lprime}(i) is exactly the new
$G(u)$: the previous layers split into those anchored $\le X-1$ and the
$K$ anchored in $[X,y]$. At termination all anchors lie in $[b]$, so
the layer count is $G(b)\le D(b)$.
\end{proof}

\begin{proof}[Proof of \Cref{lemma:Lprime}]
By \Cref{prop:det-Lprime}, deterministically $L \le D(b)$. Now fix a
bin $\tau$ and an interval $J$ of length $g$. Since the starting points
are independent and uniform over $[m]$, exactly $g$ of the $m$
positions map to bin $\tau$ with block in $J$, so
$n_\tau(J)\sim\Bin(n, g/m)$. By the definition of $C(g)$ and a union
bound over the $N$ bins and fewer than $b^2$ intervals, except with
probability $2^{-\lambda}$,
\[
  M(J)\le C(|J|)
\]
holds simultaneously for every $J$. Set
$\rho := \max_{1\le g\le b} \frac{C(g)+1}{g+\zeta-1}$, so that
$C(g) \le \rho\,(g+\zeta-1)-1$ for every $g$ and
$L' = \lceil\rho\,(b+\zeta-1)\rceil$. On this event, for any admissible
family $J_1,\ldots,J_K$, the extensions are pairwise-disjoint
subintervals of the anchor range $[2-\zeta,\,b]$, whose length is
$b+\zeta-1$, so
\[
  \sum_k M(J_k)
  \le \sum_k C(|J_k|)
  \le \rho\sum_k \bigl(|J_k|+\zeta-1\bigr) - K
  \le \rho(b+\zeta-1)
  \le L'.
\]
Taking the maximum over families gives $D(b)\le L'$, and combining the
two bounds, $\Pr[L>L']\le 2^{-\lambda}$.

For the last claim, assume $\zeta \ge \min(\lambda,b)$ and
$\lambda \ge \log_2(Nb^2)$. Let
$\lambda' := \lambda+\log_2(Nb^2) \le 2\lambda$ denote the tail level
of the cap $C(g)$, and $\mu_g := ng/m \le g$ its mean. A Bernstein tail
bound gives
$C(g) \le \mu_g + \sqrt{2\ln 2\cdot\lambda'\mu_g} + O(\lambda')$, and
with $\sqrt{\lambda'\mu_g}\le(\lambda'+\mu_g)/2$ this is
$C(g) = O(g+\lambda)$ for every $g$. If $\zeta \ge b$, then
$\tfrac{b+\zeta-1}{g+\zeta-1} \le 2$, so every term of $L'$ is at most
$2\,(C(b)+1) = O(n/N+\lambda)$, using that $C$ is nondecreasing.
Otherwise $\lambda \le \zeta \le b$, and every term is at most
\[
  O(g+\lambda)\cdot\frac{b+\zeta}{g+\zeta}
  \;=\; O(b+\zeta) + O\Bigl(\lambda\cdot\frac{b+\zeta}{\zeta}\Bigr)
  \;=\; O(b+\lambda),
\]
using $g/(g+\zeta)\le1$ and $\lambda/\zeta\le1$. Since $b=m/N=O(n/N)$,
the maximum over $g$ is $O(n/N+\lambda)$.
\end{proof}

%%%%%%%%%%%%%%%%%%%%%%%%%%%%%%%%%%%%%%%%%%%%%%%%%%%%%%%%%%%
\section{Full-field Random Band Widths}\label{sec:OKVS experiment}

\paragrap{Methodology for random-band parameter choices of \cite{BPSY23}}
Given an input size $n$ and an expansion ratio $\varepsilon$, the band width $w$ is chosen so that the induced RB matrix $\mat M \in \F^{n\times m}$ is full rank with probability at least $1 - 2^{-\lambda}$. Typically $\lambda$ is taken around $40$, but it is infeasible to tune parameters while verifying a failure probability of $2^{-40}$. To address this, \cite{BPSY23} heuristically established the linear relation between
\begin{equation}\label{eq:band_width}
    \lambda = a_{\varepsilon}\cdot w + g_{\varepsilon,n},
\end{equation}
whose core assumption is that the slope $a_\varepsilon$ depends only on $\varepsilon$, but not on an input size $n$. Then, for a fixed $\varepsilon$, one measures the slope and intercept by varying $n$ and $w$ in a regime where the failure probability is feasibly measurable (e.g., $\lambda \le 15$), and extrapolates the fitted relation to determine the value of $w$ for $\lambda = 40$.

\paragrap{Band-width comparison between binary-band and $\Z_p$-band}
For our random spaced bands over \(\Z_p\), we choose the band width \(w\) using the same full rank criterion as in \cite{BPSY23}: for each expansion ratio \(\varepsilon\), \(w\) is selected so that the induced matrix has full rank with probability at least \(1-2^{-\lambda}\).
\Cref{fig:comparisonband} compares the band width $w$ required to achieve the target failure probability in the binary-band setting and in our $\Z_p$-band setting, and shows that the $\Z_p$-band choice consistently requires smaller $w$.
Since~\cite{BPSY23} does not report binary-band parameters for all expansion ratios $\varepsilon$ considered in our study, we re-implemented the binary-band baseline and ran the corresponding experiments over a range of $\varepsilon$.

Finally, to seek the parameters for our implementation~\Cref{sec:performance}, we follow the empirical fitting strategy of~\cite{BPSY23}: measuring failure probabilities $2^{-\lambda}$ for $\lambda<15$ in the $\Z_p$-band setting and fit them to the linear model~\cref{eq:band_width}. The results are summarized in \Cref{tab:our-best-fit-lines}
%%%%%%%%%%%%%%%%Fig%%%%%%%%%%%%%%%%%%%%%%%%%%%%

\begin{figure*}[h]   
    \centering
    % ------- 1st row -------
    \begin{subfigure}[htbp]{0.32\textwidth}
        \centering
        \includegraphics[width=\textwidth]{images/1.05comp.png}
        \caption{Failure probability for $\varepsilon = 0.05$}
        \label{fig:ratio-1.05}
    \end{subfigure}
    \hfill
    \begin{subfigure}[h]{0.32\textwidth}
        \centering
        \includegraphics[width=\textwidth]{images/1.1comp.png}
        \caption{Failure probability for $\varepsilon = 0.1$}
        \label{fig:ratio-1.1}
    \end{subfigure}
    \hfill
    \begin{subfigure}[h]{0.32\textwidth}
        \centering
        \includegraphics[width=\textwidth]{images/1.3comp.png}
        \caption{Failure probability for $\varepsilon = 0.3$}
        \label{fig:ratio-1.15}
    \end{subfigure}
    \caption{Comparison of the security parameter $\lambda$ as a function of the band width $w$ between our binary-band implementation and the $\mathbb{Z}_p$-band setting. Dashed lines denote the extrapolated linear fits from the binary-band measurements, while solid lines denote the extrapolated fits over $\mathbb{Z}_p$-band. For each $n \in \{2^{16}, 2^{18}, 2^{20}\}$, the solid curve lies to the left of the dashed curve, indicating that achieving the same target security level $\lambda$ requires a smaller band width $w$ in the $\mathbb{Z}_p$-band setting.}
    \label{fig:comparisonband}
\end{figure*}

%%%%%%%%%%%%%%%Table%%%%%%%%%%%%%%%%%%%
\begin{table}[htbp]
\caption{Best-fit lines~\eqref{eq:band_width} for $\mathbb{Z}_p$-band, and the corresponding band width $w$ for the encoding failure probability $2^{-40}$}
\label{tab:our-best-fit-lines}
\centering
\small
\setlength{\tabcolsep}{20pt}
\renewcommand{\arraystretch}{1.15}
\begin{tabular}{ccc}
\toprule
\multicolumn{3}{c}{$n=2^{20}$} \\
\cmidrule(r){1-3} 
$\varepsilon$ & Fit line & $w$ for $\lambda = 40$ \\
\cmidrule(r){1-3}
 0.1 & $0.2753w - 14.0200$ & 197 \\
 0.2 & $0.4934w - 15.1968$ & 112 \\
 0.3 & $0.6956w - 15.3069$ & 80 \\
 0.4 & $0.7929w - 14.1219$ & 69 \\
 0.5 & $0.9168w - 13.7755$ & 59 \\
 0.6 & $1.1160w - 15.3274$ & 50 \\
 0.7 & $1.3036w - 15.7176$ & 43 \\
 0.8 & $1.3341w - 15.3561$ & 42 \\
 0.9 & $1.5228w - 16.6124$ & 38 \\
 1.0 & $1.7285w - 17.4118$ & 34 \\
 1.1 & $1.9032w - 19.0024$ & 31 \\
 1.2 & $1.9764w - 19.0728$ & 30 \\
 1.3 & $1.9884w - 18.2284$ & 30 \\
 1.4 & $2.1426w - 16.9244$ & 27 \\
 1.5 & $2.1845w - 17.2011$ & 27 \\
 1.6 & $2.2886w - 18.8276$ & 26 \\
 1.7 & $2.3703w - 19.0591$ & 25 \\
 1.8 & $2.5023w - 19.1985$ & 24 \\
 1.9 & $2.6012w - 19.0890$ & 23 \\
 2.0 & $2.6760w - 17.4406$ & 22 \\
\bottomrule
\end{tabular}
\end{table}

%%%%%%%%%%%%%%%%%%%%%%%%%%%%%%%%%%%%%%%%%%%%%%%%%%%%%%%%%%
\section{Asymptotic Cost Analysis}\label{sec:asymp_cost_anal}
We analyze the protocol of \Cref{fig:our-ss-pmt-protocol}, where
$n=n_x=n_y$. Since $\mathcal{F}_{\mathsf{ss\mbox{-}peqt}}$ over $NL'$
instances of $O(\lambda)$-bit values can be realized with
$O(NL'\lambda) = O(n)$ communication and computation, it suffices to
analyze the RLWE/OKVS components. Throughout we use
$L' = O(n/N+\lambda)$ from \Cref{lemma:Lprime}; its hypothesis
$\zeta \ge \min(\lambda,b)$ and $\lambda \ge \log_2(Nb^2)$ holds for
all parameters of \Cref{tab:rsb_params}.

For the communication cost, $\receiver$ sends $b+w-1$ RLWE ciphertexts
to encrypt the OKVS encoding, where $b=m/N$; each has size $O(N)$,
giving total communication $O((b+w-1)N)=O(n+Nw)=O(n)$, where the last
equality holds because $N$ is fixed and $w=O(\lambda)$. In the other
direction, $\sender$ sends exactly $L'$ RLWE ciphertexts, so the return
communication is $O((n/N+\lambda)N)=O(n+N\lambda)=O(n)$ as well.

For the computation cost, we invoke $\mathsf{RLWE.Enc}$ exactly $b+w-1$
times. Since each encryption call takes $O(N\log N)$ unit cost, the
total encryption cost is $O((b+w-1)N\log N)=O((n+wN)\log N)=O(n)$. In
$\mathsf{RSB\mbox{-}OKVS.HomDecode}$, each of the $L\le L'$ actual
layers uses $w+\zeta$ slot-wise multiplications of $O(N\log N)$ each,
and each padded layer uses $O(1)$; since $w=O(\lambda)$ and our choices
of $\zeta$ exceed $\min(\lambda,b)$ by at most a small constant factor,
each layer costs $O(\lambda)$ operations, and the total decoding cost
is $O(n\lambda)$, which is linear in $n$ for fixed $\lambda$. Finally,
$\mathsf{RLWE.Dec}$ is invoked $L'$ times at cost $O(N\log N)$ each,
for total decryption cost $O(n)$.

%%%%%%%%%%%%%%%%%%%%%%%%%%%%%%%%%%%%
%PT26과 비교하는 내용은 삭제해도 될 듯
% \section{Communication Cost Analysis of~\cite{PT26}}
% \label{app:pt26-comm}

% % As the original paper of \cite{PT26} provided no concrete communication cost of their proposed PSU protocol, we provide it here. The main purpose is to argue the near-linear dependency on the item length $L$, visualized in \Cref{fig:pt26-item-length}.

% \paragrap{Notation}
% Let $n_b=|X_b|$ denote each party's set size, $\tau=n_0+n_1$, and
% $u=|X_0\cup X_1|\le\tau$. The Invertible Bloom Lookup Table (IBLT) is parameterized by the number of hash functions $k$, the subtable length $\ell$, and the total number of bins $m=k\ell\simeq e\tau$, where $e$ is the IBLT expansion factor. Let
% $\kappa$ denote the computational security parameter. We write $L$ for
% the bit-length of the encoded value carried in the IBLT sum field and in
% the \textsf{UnionPeel} messages, including the encoding overhead for the
% special symbol $\bot$. For exact PSU over arbitrary $a$-bit items under
% the direct one-to-one encoding assumed in~\cite{PT26}, one has $L\ge a$
% up to this encoding overhead.

% \paragrap{Per-bin payload}
% The protocol $\Pi_{\mathrm{PSU}}$ of~\cite[Fig.~5]{PT26} repeatedly
% invokes the sub-protocol $\Pi_{\mathrm{UnionPeel}}$
% of~\cite[Fig.~4]{PT26} on a public bin set $Q_{\pi}^{(t)}$. In the
% optimized version of~\cite{PT26}, the per-bin OPRF calls are replaced by
% a one-time multi-point OPRF on $X_1$ during setup. We therefore exclude
% this setup term and count only the repeated per-bin OT/reveal payload.

% For each examined bin, \textsf{UnionPeel} communicates
% \begin{equation}
% \label{eq:pt26-cbin}
% \begin{aligned}
% c_{\mathrm{bin}}(L)
% &=
% \underbrace{2L}_{\text{$1$-out-of-$2$ OT}}
% +
% \underbrace{3\max(L,\kappa)}_{\text{$1$-out-of-$3$ OT}}
% +
% \underbrace{L}_{\substack{\text{reveal}\\ P_1\to P_0}}        \\
% &=
% 3L+3\max(L,\kappa)
% \quad\text{bits}.
% \end{aligned}
% \end{equation}
% The $1$-out-of-$3$ OT term uses standard bit-string OT extension: its
% mixed-domain messages are encoded as bit strings of common length
% $\max(L,\kappa)$. We also assume that the public bin set
% $Q_{\pi}^{(t)}$ is enumerated in a fixed common order, so the bin index
% $(i,j)$ need not be retransmitted. If indices are sent explicitly, an
% additional $\lceil\log_2 k\rceil+\lceil\log_2\ell\rceil$ bits per
% examined bin should be added.

% \paragrap{Total number of bin evaluations}
% Let $q_t=|Q_\pi^{(t)}|$. From
% $Q_\pi^{(0)}=[k]\times[\ell]$ and
% $Q_\pi^{(t+1)}
% =
% \{(i,h_i(v)): i\in[k],\, v\in V_\pi^{(t)}\}$
% in~\cite[Fig.~5]{PT26}, we have $q_0=m$ and
% $q_{t+1}\le k|V_\pi^{(t)}|$. Since the recovered elements
% $\bigcup_t V_\pi^{(t)}$ form $X_0\cup X_1$ on a successful execution,
% \begin{equation}
% \label{eq:pt26-qsigma}
% Q_\Sigma
% :=
% \sum_{t\ge 0} q_t
% \le
% m+k |V_{\pi}^{(t)}|
% \le
% e\tau+k\tau
% =
% (e+k)\tau .
% \end{equation}

% \paragrap{Closed-form bound}
% Combining~\eqref{eq:pt26-cbin} and~\eqref{eq:pt26-qsigma}, the dominant
% item-length-dependent payload satisfies
% \begin{equation}
% \label{eq:pt26-payload}
% C_{\mathrm{payload}}(L)
% \le
% (e+k)\tau\bigl(3L+3\max(L,\kappa)\bigr)
% \quad\text{bits}.
% \end{equation}
% For $\kappa=128$, the per-bin payload factor is piecewise-linear:
% \begin{equation}
% \label{eq:pt26-cbin-piecewise}
% c_{\mathrm{bin}}(L)
% =
% \begin{cases}
% 3L+384, & L\le 128,\\
% 6L,     & L>128.
% \end{cases}
% \end{equation}
% Thus
% \[
% C_{\mathrm{payload}}(L)
% \le
% (e+k)\tau\cdot c_{\mathrm{bin}}(L),
% \]
% and the slope doubles at $L=\kappa$ because the $1$-out-of-$3$ OT
% payload also becomes $L$-dominated.

% Normalizing by the $64$-bit baseline used in~\cite[Table~2]{PT26}
% eliminates the dependence on $\tau$, $e$, and $k$:
% \begin{equation}
% \label{eq:pt26-ratio}
% \frac{C_{\mathrm{payload}}(L)}
%      {C_{\mathrm{payload}}(64)}
% =
% \frac{L+\max(L,128)}{192}.
% \end{equation}
% This gives the ratios $1.33\times$, $2.67\times$, and $5.33\times$ at
% $L=128$, $256$, and $512$ bits, respectively.

% %-------------------------------------------------------------------------------
% \section*{Acknowledgments}
% %-------------------------------------------------------------------------------

% The USENIX latex style is old and very tired, which is why
% there's no \textbackslash{}acks command for you to use when
% acknowledging. Sorry.

% \textbf{Do not include any acknowledgements in your submission which may deanonymize you (e.g., because of specific affiliations or grants you acknowledge)}

%-------------------------------------------------------------------------------

%%%%%%%%%%%%%%%%%%%%%%%%%%%%%%%%%%%%%%%%%%%%%%%%%%%%%%%%%%%%%%%%%%%%%%%%%%%%%%%%
\end{document}
%%%%%%%%%%%%%%%%%%%%%%%%%%%%%%%%%%%%%%%%%%%%%%%%%%%%%%%%%%%%%%%%%%%%%%%%%%%%%%%%

%%  LocalWords:  endnotes includegraphics fread ptr nobj noindent
%%  LocalWords:  pdflatex acks
